\documentclass[12pt, ukrainian]{article}
\usepackage[a4paper]{geometry}
\setlength\parindent{0mm}
\geometry{top=1.1cm,bottom=1.1cm,left=1.1cm,right=1.1cm}
\pagestyle{empty}
\usepackage[T2A]{fontenc}
\usepackage[utf8]{inputenc}
\usepackage[ukrainian]{babel}
\usepackage{graphicx}
\usepackage{caption}
\DeclareCaptionFormat{nocaption}{}
\captionsetup{format=nocaption,aboveskip=0pt,belowskip=0pt}
\usepackage[Export]{adjustbox}
\adjustboxset{max size={0.9\linewidth}{0.9\paperheight}}
\begin{document}
{{ 27.11.2025 Контрольна робота \No 2, варіант  \No \textbf { 1 }\par}}
{
%\begin {large}

{%Завдання 1-10: Що буде виведено за виконання (повідомлення інтерпретатора про помилку %позначати error)?

%Завдання 11-12 (на звороті): написати функцію для розв'язання задачі. Оцінюється економність обчислень та якість коду.

%Тематика задач: використання циклів, програмування з використанням рекурентних співвідношень,
%обробка списків та рядків.

У завданнях 1-10 вказати, що буде виведено за виконання.

Повідомлення інтерпретатора про помилку позначати ERROR

Завдання 11 та 12 на звороті.

\begin{minipage}[t]{7.5cm}

\begin{center}
\textbf{Завдання 1}
\end{center}

\begin{verbatim}
a, b, c = 8, 6, 7
def g(a, b, c=9):
    print(a, b, c, end=" ")

g(5, b=1)
print(a, b, c)
\end{verbatim}



\begin{center}
\textbf{Завдання 2}
\end{center}

\begin{verbatim}
a,b,c=6,0,3
def g(b):
    global c
    a=5
    b=3
    c=4
    return a+b+c

a,b,c=1,2,8
print(g(b),a,b,c)
\end{verbatim}



\begin{center}
\textbf{Завдання 3}
\end{center}

\begin{verbatim}
def f():
    try:
        res = int(4//2)
        return 45
    except KeyboardInterrupt: 
        return 3
    except TypeError: 
        return 6
    return res

print(f())
\end{verbatim}



\begin{center}
\textbf{Завдання 4}
\end{center}

\begin{verbatim}
k=9
while k>=1:
    k+=-2
print(k)
\end{verbatim}

\end{minipage}
\vline \hspace{6mm}
\begin{minipage}[t]{10.5cm}

\begin{center}
\textbf{Завдання 5}
\end{center}

\begin{verbatim}
for a in range(9,9+3,-2):
    if a<3:
        continue
    print(a, end=' ')
    a=-7
else:
    print(a, end=' ')
print(a, end=' ')
\end{verbatim}



\begin{center}
\textbf{Завдання 6}
\end{center}

\begin{verbatim}
def f(s):
    for i in range(len(s)):
        s[i]+=1
    return

s=[1,2,3]
f(s)
print(s)
\end{verbatim}



\begin{center}
\textbf{Завдання 7}
\end{center}

\begin{verbatim}
d = ('a','b','c','d','e',0,1,2,4)
print(d[-13])
\end{verbatim}



\begin{center}
\textbf{Завдання 8}
\end{center}

\begin{verbatim}
g = '0123456789'
print(g[::-2])
\end{verbatim}



\begin{center}
\textbf{Завдання 9}
\end{center}

\begin{verbatim}
c=['0','1','2','3']
c.remove('1')
print(c)
\end{verbatim}



\begin{center}
\textbf{Завдання 10}
\end{center}
Що буде надруковано за виконання?\par
\begin{verbatim}
a=[10,11,12,13]
a[0:]=[2,3]
print(a)
\end{verbatim}

\end{minipage}

%\end {large}
}
\newpage

{{ 27.11.2025 Контрольна робота \No 2, варіант  \No \textbf { 2 }\par}}
{
%\begin {large}

{%Завдання 1-10: Що буде виведено за виконання (повідомлення інтерпретатора про помилку %позначати error)?

%Завдання 11-12 (на звороті): написати функцію для розв'язання задачі. Оцінюється економність обчислень та якість коду.

%Тематика задач: використання циклів, програмування з використанням рекурентних співвідношень,
%обробка списків та рядків.

У завданнях 1-10 вказати, що буде виведено за виконання.

Повідомлення інтерпретатора про помилку позначати ERROR

Завдання 11 та 12 на звороті.

\begin{minipage}[t]{7.5cm}

\begin{center}
\textbf{Завдання 1}
\end{center}

\begin{verbatim}
a, b, c = 7, 6, 8
def f(a, b, c):
    print(a, b, c, end=" ")

f(a=0, 3, c=4)
print(a, b, c)
\end{verbatim}



\begin{center}
\textbf{Завдання 2}
\end{center}

\begin{verbatim}
a,b,c=4,7,8
def f(b):
    a*=2
    b=4
    c=1
    return a+b+c

a,b,c=9,6,5
print(f(b),a,b,c)
\end{verbatim}



\begin{center}
\textbf{Завдання 3}
\end{center}

\begin{verbatim}
def f():
    try:
        res = int(7%0)
    except BaseException: 
        return 5
    except Exception: 
        return 1
    else: return 35
    finally: return 25
    return res

print(f())
\end{verbatim}



\begin{center}
\textbf{Завдання 4}
\end{center}

\begin{verbatim}
j=5
while j>8:
    j+=-3
print(j)
\end{verbatim}

\end{minipage}
\vline \hspace{6mm}
\begin{minipage}[t]{10.5cm}

\begin{center}
\textbf{Завдання 5}
\end{center}

\begin{verbatim}
for c in range(1,1+5,-3):
    if c<=5:
        pass
    print(c, end=' ')
    c=-8
else:
    print(c, end=' ')
print(c, end=' ')
\end{verbatim}



\begin{center}
\textbf{Завдання 6}
\end{center}

\begin{verbatim}
def f(s):
    s1=s
    for i in range(len(s)):
        s1[i]+=1
    return s1

s=[1,2,3]
f(s)
print(s)
\end{verbatim}



\begin{center}
\textbf{Завдання 7}
\end{center}

\begin{verbatim}
e = ('a','b','c','d','e',0,1)
print(e[3])
\end{verbatim}



\begin{center}
\textbf{Завдання 8}
\end{center}

\begin{verbatim}
b = ['0','1','2','3','4','5','6','7','8','9']
print(b[:-14:2])
\end{verbatim}



\begin{center}
\textbf{Завдання 9}
\end{center}

\begin{verbatim}
e=[10,11,12,13]
e.append([1,2])
print(e)
\end{verbatim}



\begin{center}
\textbf{Завдання 10}
\end{center}
Що буде надруковано за виконання?\par
\begin{verbatim}
c=['0','1','2','3']
c[-8:1]=[14,]
print(c)
\end{verbatim}

\end{minipage}

%\end {large}
}
\newpage

{{ 27.11.2025 Контрольна робота \No 2, варіант  \No \textbf { 3 }\par}}
{
%\begin {large}

{%Завдання 1-10: Що буде виведено за виконання (повідомлення інтерпретатора про помилку %позначати error)?

%Завдання 11-12 (на звороті): написати функцію для розв'язання задачі. Оцінюється економність обчислень та якість коду.

%Тематика задач: використання циклів, програмування з використанням рекурентних співвідношень,
%обробка списків та рядків.

У завданнях 1-10 вказати, що буде виведено за виконання.

Повідомлення інтерпретатора про помилку позначати ERROR

Завдання 11 та 12 на звороті.

\begin{minipage}[t]{7.5cm}

\begin{center}
\textbf{Завдання 1}
\end{center}

\begin{verbatim}
a, b, c = 9, 7, 8
def h(a, b=6, c):
    print(a, b, c, end=" ")

h(a=2, b=5, c=0)
print(a, b, c)
\end{verbatim}



\begin{center}
\textbf{Завдання 2}
\end{center}

\begin{verbatim}
a,b,c=2,3,1
def g(a):
    global c
    a+=5
    b=3
    c=5
    return a+b+c

a,b,c=0,6,9
print(g(a),a,b,c)
\end{verbatim}



\begin{center}
\textbf{Завдання 3}
\end{center}

\begin{verbatim}
def f():
    try:
        res = int("b2")
        return 43
    except ValueError: 
        return 9
    except ZeroDivisionError: 
        return 0
    else: return 33
    finally: return 23
    return res

print(f())
\end{verbatim}



\begin{center}
\textbf{Завдання 4}
\end{center}

\begin{verbatim}
n=14
while n>0:
    n+=-2
print(n)
\end{verbatim}

\end{minipage}
\vline \hspace{6mm}
\begin{minipage}[t]{10.5cm}

\begin{center}
\textbf{Завдання 5}
\end{center}

\begin{verbatim}
for d in range(0,0+4,2):
    if d>4:
        break
    print(d, end=' ')
    d=-2
else:
    print(d, end=' ')
print(d, end=' ')
\end{verbatim}



\begin{center}
\textbf{Завдання 6}
\end{center}

\begin{verbatim}
def f(n):
    if n > 9:
        f(n // 100)
    print(n % 10, end=' ')
    

f(123456)
\end{verbatim}



\begin{center}
\textbf{Завдання 7}
\end{center}

\begin{verbatim}
h = ('a','b','c','d','e',0,1,2,3)
print(h[-12])
\end{verbatim}



\begin{center}
\textbf{Завдання 8}
\end{center}

\begin{verbatim}
e = ('0','1','2','3','4','5','6','7','8','9')
print(e[15:15:-1])
\end{verbatim}



\begin{center}
\textbf{Завдання 9}
\end{center}

\begin{verbatim}
d=['0','1','2','3']
d.insert(3, '11')
print(d)
\end{verbatim}



\begin{center}
\textbf{Завдання 10}
\end{center}
Що буде надруковано за виконання?\par
\begin{verbatim}
g=[10,11,12,13]
del g[-4:-5]
print(g)
\end{verbatim}

\end{minipage}

%\end {large}
}
\newpage

{{ 27.11.2025 Контрольна робота \No 2, варіант  \No \textbf { 4 }\par}}
{
%\begin {large}

{%Завдання 1-10: Що буде виведено за виконання (повідомлення інтерпретатора про помилку %позначати error)?

%Завдання 11-12 (на звороті): написати функцію для розв'язання задачі. Оцінюється економність обчислень та якість коду.

%Тематика задач: використання циклів, програмування з використанням рекурентних співвідношень,
%обробка списків та рядків.

У завданнях 1-10 вказати, що буде виведено за виконання.

Повідомлення інтерпретатора про помилку позначати ERROR

Завдання 11 та 12 на звороті.

\begin{minipage}[t]{7.5cm}

\begin{center}
\textbf{Завдання 1}
\end{center}

\begin{verbatim}
a, b, c = 9, 8, 7
def g(a, b=9, c=6):
    print(a, b, c, end=" ")

g(1, 2, a=3)
print(a, b, c)
\end{verbatim}



\begin{center}
\textbf{Завдання 2}
\end{center}

\begin{verbatim}
a,b,c=8,7,0
def h(a):
    global c
    a=2
    b-=1
    c=3
    return a+b+c

a,b,c=1,4,3
print(h(b),a,b,c)
\end{verbatim}



\begin{center}
\textbf{Завдання 3}
\end{center}

\begin{verbatim}
def f():
    try:
        res = 8>3
    except ZeroDivisionError: 
        return 4
    except KeyboardInterrupt: 
        return 1
    return res

print(f())
\end{verbatim}



\begin{center}
\textbf{Завдання 4}
\end{center}

\begin{verbatim}
n=3
while n<=7:
    n+=2
print(n)
\end{verbatim}

\end{minipage}
\vline \hspace{6mm}
\begin{minipage}[t]{10.5cm}

\begin{center}
\textbf{Завдання 5}
\end{center}

\begin{verbatim}
for b in range(2,2+2,3):
    if b>=6:
        continue
    print(b, end=' ')
    b=0
    if b<=7:
        break
else:
    print(b, end=' ')
print(b, end=' ')
\end{verbatim}



\begin{center}
\textbf{Завдання 6}
\end{center}

\begin{verbatim}
def f(s):
    for el in s:
        el+=1
    return s

s=[1,2,3]
s=f(s)
print(s)
\end{verbatim}



\begin{center}
\textbf{Завдання 7}
\end{center}

\begin{verbatim}
f = ('a','b','c','d','e',0,1,2,3,4)
print(f[-10])
\end{verbatim}



\begin{center}
\textbf{Завдання 8}
\end{center}

\begin{verbatim}
f = [0,1,2,3,4,5,6,7,8,9]
print(f[8::3])
\end{verbatim}



\begin{center}
\textbf{Завдання 9}
\end{center}

\begin{verbatim}
a=[10,11,12,13]
print(a.pop(-4), end=' ')
print(a)
\end{verbatim}



\begin{center}
\textbf{Завдання 10}
\end{center}
Що буде надруковано за виконання?\par
\begin{verbatim}
d=['0','1','2','3']
del d[7:6]
print(d)
\end{verbatim}

\end{minipage}

%\end {large}
}
\newpage

{{ 27.11.2025 Контрольна робота \No 2, варіант  \No \textbf { 5 }\par}}
{
%\begin {large}

{%Завдання 1-10: Що буде виведено за виконання (повідомлення інтерпретатора про помилку %позначати error)?

%Завдання 11-12 (на звороті): написати функцію для розв'язання задачі. Оцінюється економність обчислень та якість коду.

%Тематика задач: використання циклів, програмування з використанням рекурентних співвідношень,
%обробка списків та рядків.

У завданнях 1-10 вказати, що буде виведено за виконання.

Повідомлення інтерпретатора про помилку позначати ERROR

Завдання 11 та 12 на звороті.

\begin{minipage}[t]{7.5cm}

\begin{center}
\textbf{Завдання 1}
\end{center}

\begin{verbatim}
a, b, c = 8, 7, 9
def f(a, b, c=6):
    print(a, b, c, end=" ")

f(4, 4)
print(a, b, c)
\end{verbatim}



\begin{center}
\textbf{Завдання 2}
\end{center}

\begin{verbatim}
a,b,c=5,2,6
def f(b):
    a-=4
    b=3
    c=5
    return a+b+c

a,b,c=3,8,9
print(f(a),a,b,c)
\end{verbatim}



\begin{center}
\textbf{Завдання 3}
\end{center}

\begin{verbatim}
def f():
    try:
        res = int("8")
        return 40
    except TypeError: 
        return 5
    except Exception: 
        return 6
    return res

print(f())
\end{verbatim}



\begin{center}
\textbf{Завдання 4}
\end{center}

\begin{verbatim}
m=15
while m>=2:
    m+=-3
print(m)
\end{verbatim}

\end{minipage}
\vline \hspace{6mm}
\begin{minipage}[t]{10.5cm}

\begin{center}
\textbf{Завдання 5}
\end{center}

\begin{verbatim}
for f in range(4,4+6,2):
    if f>=8:
        pass
    print(f, end=' ')
    f=9
else:
    print(13, end=' ')
print(f, end=' ')
\end{verbatim}



\begin{center}
\textbf{Завдання 6}
\end{center}

\begin{verbatim}
def f(s1):
    s=s1[:]
    for i in range(len(s)):
        s[i]+=1

s=[1,2,3]
f(s)
print(s)
\end{verbatim}



\begin{center}
\textbf{Завдання 7}
\end{center}

\begin{verbatim}
b = ('a','b','c','d','e',0,1,2)
print(b[10])
\end{verbatim}



\begin{center}
\textbf{Завдання 8}
\end{center}

\begin{verbatim}
c = (0,1,2,3,4,5,6,7,8,9)
print(c[-6::-3])
\end{verbatim}



\begin{center}
\textbf{Завдання 9}
\end{center}

\begin{verbatim}
b=[10,11,12,13]
b.extend(17)
print(b)
\end{verbatim}



\begin{center}
\textbf{Завдання 10}
\end{center}
Що буде надруковано за виконання?\par
\begin{verbatim}
b=[10,11,12,13]
b[-3:]=[12]
print(b)
\end{verbatim}

\end{minipage}

%\end {large}
}
\newpage

{{ 27.11.2025 Контрольна робота \No 2, варіант  \No \textbf { 6 }\par}}
{
%\begin {large}

{%Завдання 1-10: Що буде виведено за виконання (повідомлення інтерпретатора про помилку %позначати error)?

%Завдання 11-12 (на звороті): написати функцію для розв'язання задачі. Оцінюється економність обчислень та якість коду.

%Тематика задач: використання циклів, програмування з використанням рекурентних співвідношень,
%обробка списків та рядків.

У завданнях 1-10 вказати, що буде виведено за виконання.

Повідомлення інтерпретатора про помилку позначати ERROR

Завдання 11 та 12 на звороті.

\begin{minipage}[t]{7.5cm}

\begin{center}
\textbf{Завдання 1}
\end{center}

\begin{verbatim}
a, b, c = 7, 6, 9
def h(a, b, c):
    print(a, b, c, end=" ")

h(3, c=2, b=0)
print(a, b, c)
\end{verbatim}



\begin{center}
\textbf{Завдання 2}
\end{center}

\begin{verbatim}
a,b,c=1,0,2
def h(a):
    a*=2
    b=1
    c=4
    return a+b+c

a,b,c=7,5,4
print(h(a),a,b,c)
\end{verbatim}



\begin{center}
\textbf{Завдання 3}
\end{center}

\begin{verbatim}
def f():
    try:
        res = 9<=8
    except BaseException: 
        return 0
    except ValueError: 
        return 7
    else: return 32
    finally: return 20
    return res

print(f())
\end{verbatim}



\begin{center}
\textbf{Завдання 4}
\end{center}

\begin{verbatim}
i=8
while i>3:
    i+=-2
print(i)
\end{verbatim}

\end{minipage}
\vline \hspace{6mm}
\begin{minipage}[t]{10.5cm}

\begin{center}
\textbf{Завдання 5}
\end{center}

\begin{verbatim}
for e in range(-7,-7+6,-2):
    if e<=7:
        break
    print(e, end=' ')
    e=2
else:
    print(e, end=' ')
print(e, end=' ')
\end{verbatim}



\begin{center}
\textbf{Завдання 6}
\end{center}

\begin{verbatim}
def f(n):
    print(n % 10, end=' ')
    if n > 9:
        f(n // 10)
    

f(123456)
\end{verbatim}



\begin{center}
\textbf{Завдання 7}
\end{center}

\begin{verbatim}
c = ('a','b','c','d','e',0,1,2,4)
print(c[-5])
\end{verbatim}



\begin{center}
\textbf{Завдання 8}
\end{center}

\begin{verbatim}
a = ('0','1','2','3','4','5','6','7','8','9')
print(a[:-6:-2])
\end{verbatim}



\begin{center}
\textbf{Завдання 9}
\end{center}

\begin{verbatim}
h=['0','1','2','3']
h+='10'
print(h)
\end{verbatim}



\begin{center}
\textbf{Завдання 10}
\end{center}
Що буде надруковано за виконання?\par
\begin{verbatim}
f=['0','1','2','3']
f[:-2]=(15,)
print(f)
\end{verbatim}

\end{minipage}

%\end {large}
}
\newpage

{{ 27.11.2025 Контрольна робота \No 2, варіант  \No \textbf { 7 }\par}}
{
%\begin {large}

{%Завдання 1-10: Що буде виведено за виконання (повідомлення інтерпретатора про помилку %позначати error)?

%Завдання 11-12 (на звороті): написати функцію для розв'язання задачі. Оцінюється економність обчислень та якість коду.

%Тематика задач: використання циклів, програмування з використанням рекурентних співвідношень,
%обробка списків та рядків.

У завданнях 1-10 вказати, що буде виведено за виконання.

Повідомлення інтерпретатора про помилку позначати ERROR

Завдання 11 та 12 на звороті.

\begin{minipage}[t]{7.5cm}

\begin{center}
\textbf{Завдання 1}
\end{center}

\begin{verbatim}
a, b, c = 8, 6, 9
def h(a, b=8, c=7):
    print(a, b, c, end=" ")

h(b=1, c=5, 0)
print(a, b, c)
\end{verbatim}



\begin{center}
\textbf{Завдання 2}
\end{center}

\begin{verbatim}
a,b,c=6,8,3
def g(b):
    global c
    a=4
    b+=1
    c=3
    return a+b+c

a,b,c=0,5,4
print(g(b),a,b,c)
\end{verbatim}



\begin{center}
\textbf{Завдання 3}
\end{center}

\begin{verbatim}
def f():
    try:
        res = int("2")
        return 42
    except BaseException: 
        return 3
    except Exception: 
        return 2
    finally: return 21
    return res

print(f())
\end{verbatim}



\begin{center}
\textbf{Завдання 4}
\end{center}

\begin{verbatim}
t=13
while t>=6:
    t+=-3
print(t)
\end{verbatim}

\end{minipage}
\vline \hspace{6mm}
\begin{minipage}[t]{10.5cm}

\begin{center}
\textbf{Завдання 5}
\end{center}

\begin{verbatim}
for d in range(8,8+3,-3):
    if d<6:
        continue
    print(d, end=' ')
    d=4
else:
    print(d, end=' ')
print(d, end=' ')
\end{verbatim}



\begin{center}
\textbf{Завдання 6}
\end{center}

\begin{verbatim}
def f(s):
    for el in s:
        el+=1
    return s

s=[1,2,3]
s=f(s)
print(s)
\end{verbatim}



\begin{center}
\textbf{Завдання 7}
\end{center}

\begin{verbatim}
a = ('a','b','c','d','e',0,1,2,3,4)
print(a[0])
\end{verbatim}



\begin{center}
\textbf{Завдання 8}
\end{center}

\begin{verbatim}
h = (0,1,2,3,4,5,6,7,8,9)
print(h[0:-12:-3])
\end{verbatim}



\begin{center}
\textbf{Завдання 9}
\end{center}

\begin{verbatim}
g=['0','1','2','3']
print(g.index('3'), end=' ')
\end{verbatim}



\begin{center}
\textbf{Завдання 10}
\end{center}
Що буде надруковано за виконання?\par
\begin{verbatim}
e=['0','1','2','3']
del e[:-3]
print(e)
\end{verbatim}

\end{minipage}

%\end {large}
}
\newpage

{{ 27.11.2025 Контрольна робота \No 2, варіант  \No \textbf { 8 }\par}}
{
%\begin {large}

{%Завдання 1-10: Що буде виведено за виконання (повідомлення інтерпретатора про помилку %позначати error)?

%Завдання 11-12 (на звороті): написати функцію для розв'язання задачі. Оцінюється економність обчислень та якість коду.

%Тематика задач: використання циклів, програмування з використанням рекурентних співвідношень,
%обробка списків та рядків.

У завданнях 1-10 вказати, що буде виведено за виконання.

Повідомлення інтерпретатора про помилку позначати ERROR

Завдання 11 та 12 на звороті.

\begin{minipage}[t]{7.5cm}

\begin{center}
\textbf{Завдання 1}
\end{center}

\begin{verbatim}
a, b, c = 9, 6, 8
def f(a, b=7, c):
    print(a, b, c, end=" ")

f(5, 3, 4)
print(a, b, c)
\end{verbatim}



\begin{center}
\textbf{Завдання 2}
\end{center}

\begin{verbatim}
a,b,c=5,7,9
def h(a):
    a=2
    b=1
    c=5
    return a+b+c

a,b,c=4,8,7
print(h(a),a,b,c)
\end{verbatim}



\begin{center}
\textbf{Завдання 3}
\end{center}

\begin{verbatim}
def f():
    try:
        res = int("c1")
    except KeyboardInterrupt: 
        return 4
    except TypeError: 
        return 0
    else: return 31
    return res

print(f())
\end{verbatim}



\begin{center}
\textbf{Завдання 4}
\end{center}

\begin{verbatim}
j=2
while j<5:
    j+=1
print(j)
\end{verbatim}

\end{minipage}
\vline \hspace{6mm}
\begin{minipage}[t]{10.5cm}

\begin{center}
\textbf{Завдання 5}
\end{center}

\begin{verbatim}
for c in range(5,5+4,3):
    if c>7:
        continue
    print(c, end=' ')
    c=-5
else:
    print(c, end=' ')
print(c, end=' ')
\end{verbatim}



\begin{center}
\textbf{Завдання 6}
\end{center}

\begin{verbatim}
def f(n):
    print(n % 10, end=' ')
    if n > 9:
           f(n // 100)
    

f(123456)
\end{verbatim}



\begin{center}
\textbf{Завдання 7}
\end{center}

\begin{verbatim}
g = ('a','b','c','d','e',0,1)
print(g[-1])
\end{verbatim}



\begin{center}
\textbf{Завдання 8}
\end{center}

\begin{verbatim}
d = '0123456789'
print(d[-15:-10:2])
\end{verbatim}



\begin{center}
\textbf{Завдання 9}
\end{center}

\begin{verbatim}
f=[10,11,12,13]
print(f.pop(2), end=' ')
print(f)
\end{verbatim}



\begin{center}
\textbf{Завдання 10}
\end{center}
Що буде надруковано за виконання?\par
\begin{verbatim}
h=[10,11,12,13]
h[8:2]='13'
print(h)
\end{verbatim}

\end{minipage}

%\end {large}
}
\newpage

{{ 27.11.2025 Контрольна робота \No 2, варіант  \No \textbf { 9 }\par}}
{
%\begin {large}

{%Завдання 1-10: Що буде виведено за виконання (повідомлення інтерпретатора про помилку %позначати error)?

%Завдання 11-12 (на звороті): написати функцію для розв'язання задачі. Оцінюється економність обчислень та якість коду.

%Тематика задач: використання циклів, програмування з використанням рекурентних співвідношень,
%обробка списків та рядків.

У завданнях 1-10 вказати, що буде виведено за виконання.

Повідомлення інтерпретатора про помилку позначати ERROR

Завдання 11 та 12 на звороті.

\begin{minipage}[t]{7.5cm}

\begin{center}
\textbf{Завдання 1}
\end{center}

\begin{verbatim}
a, b, c = 7, 9, 8
def g(a, b, c=6):
    print(a, b, c, end=" ")

g(a=1, 2, b=4)
print(a, b, c)
\end{verbatim}



\begin{center}
\textbf{Завдання 2}
\end{center}

\begin{verbatim}
a,b,c=1,2,9
def h(b):
    a-=2
    b=5
    c=1
    return a+b+c

a,b,c=7,1,6
print(h(b),a,b,c)
\end{verbatim}



\begin{center}
\textbf{Завдання 3}
\end{center}

\begin{verbatim}
def f():
    try:
        res = int(8/1)
        return 44
    except ValueError: 
        return 6
    except ZeroDivisionError: 
        return 7
    finally: return 22
    return res

print(f())
\end{verbatim}



\begin{center}
\textbf{Завдання 4}
\end{center}

\begin{verbatim}
m=1
while m<=13:
    m+=1
print(m)
\end{verbatim}

\end{minipage}
\vline \hspace{6mm}
\begin{minipage}[t]{10.5cm}

\begin{center}
\textbf{Завдання 5}
\end{center}

\begin{verbatim}
for e in range(-6,-6+2,2):
    if e>=3:
        pass
    print(e, end=' ')
    e=6
else:
    print(e, end=' ')
print(e, end=' ')
\end{verbatim}



\begin{center}
\textbf{Завдання 6}
\end{center}

\begin{verbatim}
def f(s):
    for i in range(len(s)):
        s[i]+=1
    return

s=[1,2,3]
f(s)
print(s)
\end{verbatim}



\begin{center}
\textbf{Завдання 7}
\end{center}

\begin{verbatim}
d = ('a','b','c','d','e',0,1,2)
print(d[-8])
\end{verbatim}



\begin{center}
\textbf{Завдання 8}
\end{center}

\begin{verbatim}
f = ['0','1','2','3','4','5','6','7','8','9']
print(f[-7:11:3])
\end{verbatim}



\begin{center}
\textbf{Завдання 9}
\end{center}

\begin{verbatim}
e=['0','1','2','3']
e.insert(-5, '12')
print(e)
\end{verbatim}



\begin{center}
\textbf{Завдання 10}
\end{center}
Що буде надруковано за виконання?\par
\begin{verbatim}
b=[10,11,12,13]
b[-5:-6]=()
print(b)
\end{verbatim}

\end{minipage}

%\end {large}
}
\newpage

{{ 27.11.2025 Контрольна робота \No 2, варіант  \No \textbf { 10 }\par}}
{
%\begin {large}

{%Завдання 1-10: Що буде виведено за виконання (повідомлення інтерпретатора про помилку %позначати error)?

%Завдання 11-12 (на звороті): написати функцію для розв'язання задачі. Оцінюється економність обчислень та якість коду.

%Тематика задач: використання циклів, програмування з використанням рекурентних співвідношень,
%обробка списків та рядків.

У завданнях 1-10 вказати, що буде виведено за виконання.

Повідомлення інтерпретатора про помилку позначати ERROR

Завдання 11 та 12 на звороті.

\begin{minipage}[t]{7.5cm}

\begin{center}
\textbf{Завдання 1}
\end{center}

\begin{verbatim}
a, b, c = 9, 7, 8
def f(a, b=6, c=6):
    print(a, b, c, end=" ")

f(5, 3)
print(a, b, c)
\end{verbatim}



\begin{center}
\textbf{Завдання 2}
\end{center}

\begin{verbatim}
a,b,c=5,4,8
def g(a):
    global c
    a*=3
    b=2
    c=5
    return a+b+c

a,b,c=2,0,9
print(g(a),a,b,c)
\end{verbatim}



\begin{center}
\textbf{Завдання 3}
\end{center}

\begin{verbatim}
def f():
    try:
        res = int(3%0.0)
    except KeyboardInterrupt: 
        return 9
    except Exception: 
        return 5
    else: return 30
    return res

print(f())
\end{verbatim}



\begin{center}
\textbf{Завдання 4}
\end{center}

\begin{verbatim}
i=7
while i<8:
    i+=2
print(i)
\end{verbatim}

\end{minipage}
\vline \hspace{6mm}
\begin{minipage}[t]{10.5cm}

\begin{center}
\textbf{Завдання 5}
\end{center}

\begin{verbatim}
for f in range(-8,-8+5,-3):
    if f<=5:
        break
    print(f, end=' ')
    f=5
else:
    print(f, end=' ')
print(f, end=' ')
\end{verbatim}



\begin{center}
\textbf{Завдання 6}
\end{center}

\begin{verbatim}
def f(s):
    s1=s
    for i in range(len(s)):
        s1[i]+=1
    return s1

s=[1,2,3]
f(s)
print(s)
\end{verbatim}



\begin{center}
\textbf{Завдання 7}
\end{center}

\begin{verbatim}
h = ('a','b','c','d','e',0,1,2,3)
print(h[-6])
\end{verbatim}



\begin{center}
\textbf{Завдання 8}
\end{center}

\begin{verbatim}
g = [0,1,2,3,4,5,6,7,8,9]
print(g[10:1:-1])
\end{verbatim}



\begin{center}
\textbf{Завдання 9}
\end{center}

\begin{verbatim}
h=[10,11,12,13]
h.extend('13')
print(h)
\end{verbatim}



\begin{center}
\textbf{Завдання 10}
\end{center}
Що буде надруковано за виконання?\par
\begin{verbatim}
h=['0','1','2','3']
del h[-2]
print(h)
\end{verbatim}

\end{minipage}

%\end {large}
}
\newpage

{{ 27.11.2025 Контрольна робота \No 2, варіант  \No \textbf { 11 }\par}}
{
%\begin {large}

{%Завдання 1-10: Що буде виведено за виконання (повідомлення інтерпретатора про помилку %позначати error)?

%Завдання 11-12 (на звороті): написати функцію для розв'язання задачі. Оцінюється економність обчислень та якість коду.

%Тематика задач: використання циклів, програмування з використанням рекурентних співвідношень,
%обробка списків та рядків.

У завданнях 1-10 вказати, що буде виведено за виконання.

Повідомлення інтерпретатора про помилку позначати ERROR

Завдання 11 та 12 на звороті.

\begin{minipage}[t]{7.5cm}

\begin{center}
\textbf{Завдання 1}
\end{center}

\begin{verbatim}
a, b, c = 9, 8, 7
def h(a, b, c):
    print(a, b, c, end=" ")

h(1, 0, c=2)
print(a, b, c)
\end{verbatim}



\begin{center}
\textbf{Завдання 2}
\end{center}

\begin{verbatim}
a,b,c=7,3,5
def f(a):
    a=4
    b+=5
    c=1
    return a+b+c

a,b,c=2,1,9
print(f(a),a,b,c)
\end{verbatim}



\begin{center}
\textbf{Завдання 3}
\end{center}

\begin{verbatim}
def f():
    try:
        res = int("a4")
        return 41
    except ZeroDivisionError: 
        return 6
    except TypeError: 
        return 9
    else: return 34
    return res

print(f())
\end{verbatim}



\begin{center}
\textbf{Завдання 4}
\end{center}

\begin{verbatim}
t=7
while t>9:
    t+=-2
print(t)
\end{verbatim}

\end{minipage}
\vline \hspace{6mm}
\begin{minipage}[t]{10.5cm}

\begin{center}
\textbf{Завдання 5}
\end{center}

\begin{verbatim}
for a in range(3,3+2,-2):
    if a>8:
        continue
    print(a, end=' ')
    a=-1
else:
    print(13, end=' ')
print(a, end=' ')
\end{verbatim}



\begin{center}
\textbf{Завдання 6}
\end{center}

\begin{verbatim}
def f(s1):
    s=s1[:]
    for i in range(len(s)):
        s[i]+=1

s=[1,2,3]
f(s)
print(s)
\end{verbatim}



\begin{center}
\textbf{Завдання 7}
\end{center}

\begin{verbatim}
e = ('a','b','c','d','e',0,1,2,3)
print(e[-12])
\end{verbatim}



\begin{center}
\textbf{Завдання 8}
\end{center}

\begin{verbatim}
c = ['0','1','2','3','4','5','6','7','8','9']
print(c[:4:3])
\end{verbatim}



\begin{center}
\textbf{Завдання 9}
\end{center}

\begin{verbatim}
b=[10,11,12,13]
b.remove(13)
print(b)
\end{verbatim}



\begin{center}
\textbf{Завдання 10}
\end{center}
Що буде надруковано за виконання?\par
\begin{verbatim}
g=[10,11,12,13]
g[:8]='11'
print(g)
\end{verbatim}

\end{minipage}

%\end {large}
}
\newpage

{{ 27.11.2025 Контрольна робота \No 2, варіант  \No \textbf { 12 }\par}}
{
%\begin {large}

{%Завдання 1-10: Що буде виведено за виконання (повідомлення інтерпретатора про помилку %позначати error)?

%Завдання 11-12 (на звороті): написати функцію для розв'язання задачі. Оцінюється економність обчислень та якість коду.

%Тематика задач: використання циклів, програмування з використанням рекурентних співвідношень,
%обробка списків та рядків.

У завданнях 1-10 вказати, що буде виведено за виконання.

Повідомлення інтерпретатора про помилку позначати ERROR

Завдання 11 та 12 на звороті.

\begin{minipage}[t]{7.5cm}

\begin{center}
\textbf{Завдання 1}
\end{center}

\begin{verbatim}
a, b, c = 7, 9, 6
def g(a, b=8, c):
    print(a, b, c, end=" ")

g(a=1, c=2, a=3)
print(a, b, c)
\end{verbatim}



\begin{center}
\textbf{Завдання 2}
\end{center}

\begin{verbatim}
a,b,c=7,3,6
def h(b):
    global c
    a+=2
    b=4
    c=3
    return a+b+c

a,b,c=4,0,8
print(h(b),a,b,c)
\end{verbatim}



\begin{center}
\textbf{Завдання 3}
\end{center}

\begin{verbatim}
def f():
    try:
        res = int("0")
    except BaseException: 
        return 2
    except ValueError: 
        return 1
    finally: return 24
    return res

print(f())
\end{verbatim}



\begin{center}
\textbf{Завдання 4}
\end{center}

\begin{verbatim}
t=5
while t<11:
    t+=1
print(t)
\end{verbatim}

\end{minipage}
\vline \hspace{6mm}
\begin{minipage}[t]{10.5cm}

\begin{center}
\textbf{Завдання 5}
\end{center}

\begin{verbatim}
for b in range(-9,-9+5,3):
    if b<4:
        continue
    print(b, end=' ')
    b=-9
    if b>=5:
        break
else:
    print(b, end=' ')
print(b, end=' ')
\end{verbatim}



\begin{center}
\textbf{Завдання 6}
\end{center}

\begin{verbatim}
def f(s):
    for el in s:
        el+=1
    return s

s=[1,2,3]
s=f(s)
print(s)
\end{verbatim}



\begin{center}
\textbf{Завдання 7}
\end{center}

\begin{verbatim}
c = ('a','b','c','d','e',0,1)
print(c[5])
\end{verbatim}



\begin{center}
\textbf{Завдання 8}
\end{center}

\begin{verbatim}
a = [0,1,2,3,4,5,6,7,8,9]
print(a[-11:8:2])
\end{verbatim}



\begin{center}
\textbf{Завдання 9}
\end{center}

\begin{verbatim}
f=['0','1','2','3']
f+=[2,3]
print(f)
\end{verbatim}



\begin{center}
\textbf{Завдання 10}
\end{center}
Що буде надруковано за виконання?\par
\begin{verbatim}
c=['0','1','2','3']
c[:-7]=[1,2]
print(c)
\end{verbatim}

\end{minipage}

%\end {large}
}
\newpage

{{ 27.11.2025 Контрольна робота \No 2, варіант  \No \textbf { 13 }\par}}
{
%\begin {large}

{%Завдання 1-10: Що буде виведено за виконання (повідомлення інтерпретатора про помилку %позначати error)?

%Завдання 11-12 (на звороті): написати функцію для розв'язання задачі. Оцінюється економність обчислень та якість коду.

%Тематика задач: використання циклів, програмування з використанням рекурентних співвідношень,
%обробка списків та рядків.

У завданнях 1-10 вказати, що буде виведено за виконання.

Повідомлення інтерпретатора про помилку позначати ERROR

Завдання 11 та 12 на звороті.

\begin{minipage}[t]{7.5cm}

\begin{center}
\textbf{Завдання 1}
\end{center}

\begin{verbatim}
a, b, c = 9, 8, 7
def f(a, b=6, c=6):
    print(a, b, c, end=" ")

f(b=5, c=0, 4)
print(a, b, c)
\end{verbatim}



\begin{center}
\textbf{Завдання 2}
\end{center}

\begin{verbatim}
a,b,c=0,4,9
def g(b):
    a-=2
    b=3
    c=4
    return a+b+c

a,b,c=7,1,8
print(g(a),a,b,c)
\end{verbatim}



\begin{center}
\textbf{Завдання 3}
\end{center}

\begin{verbatim}
def f():
    try:
        res = int(7/2)
    except ValueError: 
        return 5
    except ZeroDivisionError: 
        return 3
    finally: return 23
    return res

print(f())
\end{verbatim}



\begin{center}
\textbf{Завдання 4}
\end{center}

\begin{verbatim}
j=10
while j>=7:
    j+=-3
print(j)
\end{verbatim}

\end{minipage}
\vline \hspace{6mm}
\begin{minipage}[t]{10.5cm}

\begin{center}
\textbf{Завдання 5}
\end{center}

\begin{verbatim}
for b in range(6,6+4,-2):
    if b<=8:
        continue
    print(b, end=' ')
    b=8
else:
    print(13, end=' ')
print(b, end=' ')
\end{verbatim}



\begin{center}
\textbf{Завдання 6}
\end{center}

\begin{verbatim}
def f(s):
    s1=s
    for i in range(len(s)):
        s1[i]+=1
    return s1

s=[1,2,3]
f(s)
print(s)
\end{verbatim}



\begin{center}
\textbf{Завдання 7}
\end{center}

\begin{verbatim}
a = ('a','b','c','d','e',0,1,2,3,4)
print(a[-11])
\end{verbatim}



\begin{center}
\textbf{Завдання 8}
\end{center}

\begin{verbatim}
e = (0,1,2,3,4,5,6,7,8,9)
print(e[7:9:-1])
\end{verbatim}



\begin{center}
\textbf{Завдання 9}
\end{center}

\begin{verbatim}
g=['0','1','2','3']
g.append([1,0])
print(g)
\end{verbatim}



\begin{center}
\textbf{Завдання 10}
\end{center}
Що буде надруковано за виконання?\par
\begin{verbatim}
e=['0','1','2','3']
e[5]=(13)
print(e)
\end{verbatim}

\end{minipage}

%\end {large}
}
\newpage

{{ 27.11.2025 Контрольна робота \No 2, варіант  \No \textbf { 14 }\par}}
{
%\begin {large}

{%Завдання 1-10: Що буде виведено за виконання (повідомлення інтерпретатора про помилку %позначати error)?

%Завдання 11-12 (на звороті): написати функцію для розв'язання задачі. Оцінюється економність обчислень та якість коду.

%Тематика задач: використання циклів, програмування з використанням рекурентних співвідношень,
%обробка списків та рядків.

У завданнях 1-10 вказати, що буде виведено за виконання.

Повідомлення інтерпретатора про помилку позначати ERROR

Завдання 11 та 12 на звороті.

\begin{minipage}[t]{7.5cm}

\begin{center}
\textbf{Завдання 1}
\end{center}

\begin{verbatim}
a, b, c = 8, 7, 9
def g(a, b=9, c):
    print(a, b, c, end=" ")

g(3, 0, 4)
print(a, b, c)
\end{verbatim}



\begin{center}
\textbf{Завдання 2}
\end{center}

\begin{verbatim}
a,b,c=0,3,4
def f(b):
    a=4
    b=1
    c=3
    return a+b+c

a,b,c=5,6,1
print(f(a),a,b,c)
\end{verbatim}



\begin{center}
\textbf{Завдання 3}
\end{center}

\begin{verbatim}
def f():
    try:
        res = 8>=5
        return 43
    except KeyboardInterrupt: 
        return 9
    except TypeError: 
        return 8
    else: return 31
    return res

print(f())
\end{verbatim}



\begin{center}
\textbf{Завдання 4}
\end{center}

\begin{verbatim}
i=12
while i>=4:
    i+=-2
print(i)
\end{verbatim}

\end{minipage}
\vline \hspace{6mm}
\begin{minipage}[t]{10.5cm}

\begin{center}
\textbf{Завдання 5}
\end{center}

\begin{verbatim}
for e in range(-2,-2+6,3):
    if e>4:
        break
    print(e, end=' ')
    e=-6
else:
    print(e, end=' ')
print(e, end=' ')
\end{verbatim}



\begin{center}
\textbf{Завдання 6}
\end{center}

\begin{verbatim}
def f(s):
    for i in range(len(s)):
        s[i]+=1
    return

s=[1,2,3]
f(s)
print(s)
\end{verbatim}



\begin{center}
\textbf{Завдання 7}
\end{center}

\begin{verbatim}
f = ('a','b','c','d','e',0,1,2,4)
print(f[12])
\end{verbatim}



\begin{center}
\textbf{Завдання 8}
\end{center}

\begin{verbatim}
d = ('0','1','2','3','4','5','6','7','8','9')
print(d[-8:7:-3])
\end{verbatim}



\begin{center}
\textbf{Завдання 9}
\end{center}

\begin{verbatim}
a=[10,11,12,13]
print(a.index(15), end=' ')
\end{verbatim}



\begin{center}
\textbf{Завдання 10}
\end{center}
Що буде надруковано за виконання?\par
\begin{verbatim}
d=[10,11,12,13]
del d[3:3]
print(d)
\end{verbatim}

\end{minipage}

%\end {large}
}
\newpage

{{ 27.11.2025 Контрольна робота \No 2, варіант  \No \textbf { 15 }\par}}
{
%\begin {large}

{%Завдання 1-10: Що буде виведено за виконання (повідомлення інтерпретатора про помилку %позначати error)?

%Завдання 11-12 (на звороті): написати функцію для розв'язання задачі. Оцінюється економність обчислень та якість коду.

%Тематика задач: використання циклів, програмування з використанням рекурентних співвідношень,
%обробка списків та рядків.

У завданнях 1-10 вказати, що буде виведено за виконання.

Повідомлення інтерпретатора про помилку позначати ERROR

Завдання 11 та 12 на звороті.

\begin{minipage}[t]{7.5cm}

\begin{center}
\textbf{Завдання 1}
\end{center}

\begin{verbatim}
a, b, c = 7, 8, 6
def h(a, b, c):
    print(a, b, c, end=" ")

h(5, c=1, b=2)
print(a, b, c)
\end{verbatim}



\begin{center}
\textbf{Завдання 2}
\end{center}

\begin{verbatim}
a,b,c=9,2,4
def h(a):
    global c
    a=2
    b=3
    c=2
    return a+b+c

a,b,c=3,1,5
print(h(a),a,b,c)
\end{verbatim}



\begin{center}
\textbf{Завдання 3}
\end{center}

\begin{verbatim}
def f():
    try:
        res = int(3//0)
    except Exception: 
        return 4
    except BaseException: 
        return 1
    else: return 32
    return res

print(f())
\end{verbatim}



\begin{center}
\textbf{Завдання 4}
\end{center}

\begin{verbatim}
k=9
while k<=6:
    k+=2
print(k)
\end{verbatim}

\end{minipage}
\vline \hspace{6mm}
\begin{minipage}[t]{10.5cm}

\begin{center}
\textbf{Завдання 5}
\end{center}

\begin{verbatim}
for a in range(10,10+3,-3):
    if a>=3:
        continue
    print(a, end=' ')
    a=3
else:
    print(a, end=' ')
print(a, end=' ')
\end{verbatim}



\begin{center}
\textbf{Завдання 6}
\end{center}

\begin{verbatim}
def f(n):
    if n > 9:
            f(n // 100)
    print(n % 10, end=' ')
    

f(987654)
\end{verbatim}



\begin{center}
\textbf{Завдання 7}
\end{center}

\begin{verbatim}
b = ('a','b','c','d','e',0,1,2)
print(b[6])
\end{verbatim}



\begin{center}
\textbf{Завдання 8}
\end{center}

\begin{verbatim}
b = '0123456789'
print(b[-1::-2])
\end{verbatim}



\begin{center}
\textbf{Завдання 9}
\end{center}

\begin{verbatim}
c=[10,11,12,13]
c.extend([4,5])
print(c)
\end{verbatim}



\begin{center}
\textbf{Завдання 10}
\end{center}
Що буде надруковано за виконання?\par
\begin{verbatim}
f=['0','1','2','3']
f[-6:]='10'
print(f)
\end{verbatim}

\end{minipage}

%\end {large}
}
\newpage

{{ 27.11.2025 Контрольна робота \No 2, варіант  \No \textbf { 16 }\par}}
{
%\begin {large}

{%Завдання 1-10: Що буде виведено за виконання (повідомлення інтерпретатора про помилку %позначати error)?

%Завдання 11-12 (на звороті): написати функцію для розв'язання задачі. Оцінюється економність обчислень та якість коду.

%Тематика задач: використання циклів, програмування з використанням рекурентних співвідношень,
%обробка списків та рядків.

У завданнях 1-10 вказати, що буде виведено за виконання.

Повідомлення інтерпретатора про помилку позначати ERROR

Завдання 11 та 12 на звороті.

\begin{minipage}[t]{7.5cm}

\begin{center}
\textbf{Завдання 1}
\end{center}

\begin{verbatim}
a, b, c = 6, 7, 8
def f(a, b, c=9):
    print(a, b, c, end=" ")

f(0, b=4)
print(a, b, c)
\end{verbatim}



\begin{center}
\textbf{Завдання 2}
\end{center}

\begin{verbatim}
a,b,c=3,6,5
def f(a):
    global c
    a*=5
    b=1
    c=3
    return a+b+c

a,b,c=2,3,7
print(f(b),a,b,c)
\end{verbatim}



\begin{center}
\textbf{Завдання 3}
\end{center}

\begin{verbatim}
def f():
    try:
        res = int(0/0.0)
        return 45
    except BaseException: 
        return 6
    except ValueError: 
        return 5
    finally: return 21
    return res

print(f())
\end{verbatim}



\begin{center}
\textbf{Завдання 4}
\end{center}

\begin{verbatim}
t=8
while t<=9:
    t+=2
print(t)
\end{verbatim}

\end{minipage}
\vline \hspace{6mm}
\begin{minipage}[t]{10.5cm}

\begin{center}
\textbf{Завдання 5}
\end{center}

\begin{verbatim}
for c in range(-10,-10+2,2):
    if c<7:
        pass
    print(c, end=' ')
    c=-3
else:
    print(c, end=' ')
print(c, end=' ')
\end{verbatim}



\begin{center}
\textbf{Завдання 6}
\end{center}

\begin{verbatim}
def f(s1):
    s=s1[:]
    for i in range(len(s)):
        s[i]+=1

s=[1,2,3]
f(s)
print(s)
\end{verbatim}



\begin{center}
\textbf{Завдання 7}
\end{center}

\begin{verbatim}
g = ('a','b','c','d','e',0,1,2,4)
print(g[-3])
\end{verbatim}



\begin{center}
\textbf{Завдання 8}
\end{center}

\begin{verbatim}
h = [0,1,2,3,4,5,6,7,8,9]
print(h[-12::-3])
\end{verbatim}



\begin{center}
\textbf{Завдання 9}
\end{center}

\begin{verbatim}
d=['0','1','2','3']
d.insert(0, [1,2])
print(d)
\end{verbatim}



\begin{center}
\textbf{Завдання 10}
\end{center}
Що буде надруковано за виконання?\par
\begin{verbatim}
a=[10,11,12,13]
a[-9:-9]=[1,0]
print(a)
\end{verbatim}

\end{minipage}

%\end {large}
}
\newpage

{{ 27.11.2025 Контрольна робота \No 2, варіант  \No \textbf { 17 }\par}}
{
%\begin {large}

{%Завдання 1-10: Що буде виведено за виконання (повідомлення інтерпретатора про помилку %позначати error)?

%Завдання 11-12 (на звороті): написати функцію для розв'язання задачі. Оцінюється економність обчислень та якість коду.

%Тематика задач: використання циклів, програмування з використанням рекурентних співвідношень,
%обробка списків та рядків.

У завданнях 1-10 вказати, що буде виведено за виконання.

Повідомлення інтерпретатора про помилку позначати ERROR

Завдання 11 та 12 на звороті.

\begin{minipage}[t]{7.5cm}

\begin{center}
\textbf{Завдання 1}
\end{center}

\begin{verbatim}
a, b, c = 8, 7, 9
def h(a, b, c=6):
    print(a, b, c, end=" ")

h(3, a=1)
print(a, b, c)
\end{verbatim}



\begin{center}
\textbf{Завдання 2}
\end{center}

\begin{verbatim}
a,b,c=1,0,5
def f(b):
    global c
    a=1
    b*=2
    c=4
    return a+b+c

a,b,c=9,8,6
print(f(a),a,b,c)
\end{verbatim}



\begin{center}
\textbf{Завдання 3}
\end{center}

\begin{verbatim}
def f():
    try:
        res = int("7")
    except ZeroDivisionError: 
        return 2
    except Exception: 
        return 5
    else: return 30
    return res

print(f())
\end{verbatim}



\begin{center}
\textbf{Завдання 4}
\end{center}

\begin{verbatim}
k=11
while k>5:
    k+=-3
print(k)
\end{verbatim}

\end{minipage}
\vline \hspace{6mm}
\begin{minipage}[t]{10.5cm}

\begin{center}
\textbf{Завдання 5}
\end{center}

\begin{verbatim}
for f in range(-3,-3+6,3):
    if f>6:
        pass
    print(f, end=' ')
    f=7
else:
    print(f, end=' ')
print(f, end=' ')
\end{verbatim}



\begin{center}
\textbf{Завдання 6}
\end{center}

\begin{verbatim}
def f(n):
    if n > 9:
        f(n // 10)
    print(n % 10, end=' ')
    

f(543216)
\end{verbatim}



\begin{center}
\textbf{Завдання 7}
\end{center}

\begin{verbatim}
d = ('a','b','c','d','e',0,1)
print(d[-8])
\end{verbatim}



\begin{center}
\textbf{Завдання 8}
\end{center}

\begin{verbatim}
b = (0,1,2,3,4,5,6,7,8,9)
print(b[-10:2:2])
\end{verbatim}



\begin{center}
\textbf{Завдання 9}
\end{center}

\begin{verbatim}
e=[10,11,12,13]
e+=[4,5]
print(e)
\end{verbatim}



\begin{center}
\textbf{Завдання 10}
\end{center}
Що буде надруковано за виконання?\par
\begin{verbatim}
b=[10,11,12,13]
b[-7:]=[4,5]
print(b)
\end{verbatim}

\end{minipage}

%\end {large}
}
\newpage

{{ 27.11.2025 Контрольна робота \No 2, варіант  \No \textbf { 18 }\par}}
{
%\begin {large}

{%Завдання 1-10: Що буде виведено за виконання (повідомлення інтерпретатора про помилку %позначати error)?

%Завдання 11-12 (на звороті): написати функцію для розв'язання задачі. Оцінюється економність обчислень та якість коду.

%Тематика задач: використання циклів, програмування з використанням рекурентних співвідношень,
%обробка списків та рядків.

У завданнях 1-10 вказати, що буде виведено за виконання.

Повідомлення інтерпретатора про помилку позначати ERROR

Завдання 11 та 12 на звороті.

\begin{minipage}[t]{7.5cm}

\begin{center}
\textbf{Завдання 1}
\end{center}

\begin{verbatim}
a, b, c = 7, 8, 9
def g(a, b, c):
    print(a, b, c, end=" ")

g(5, c=2, b=3)
print(a, b, c)
\end{verbatim}



\begin{center}
\textbf{Завдання 2}
\end{center}

\begin{verbatim}
a,b,c=2,4,6
def g(a):
    a=5
    b+=3
    c=2
    return a+b+c

a,b,c=1,7,5
print(g(b),a,b,c)
\end{verbatim}



\begin{center}
\textbf{Завдання 3}
\end{center}

\begin{verbatim}
def f():
    try:
        res = 8<4
        return 44
    except TypeError: 
        return 3
    except KeyboardInterrupt: 
        return 7
    finally: return 24
    return res

print(f())
\end{verbatim}



\begin{center}
\textbf{Завдання 4}
\end{center}

\begin{verbatim}
i=0
while i<12:
    i+=1
print(i)
\end{verbatim}

\end{minipage}
\vline \hspace{6mm}
\begin{minipage}[t]{10.5cm}

\begin{center}
\textbf{Завдання 5}
\end{center}

\begin{verbatim}
for d in range(-5,-5+4,2):
    if d<5:
        continue
    print(d, end=' ')
    d=1
else:
    print(d, end=' ')
print(d, end=' ')
\end{verbatim}



\begin{center}
\textbf{Завдання 6}
\end{center}

\begin{verbatim}
def f(s):
    for el in s:
        el+=1
    return s

s=[1,2,3]
s=f(s)
print(s)
\end{verbatim}



\begin{center}
\textbf{Завдання 7}
\end{center}

\begin{verbatim}
c = ('a','b','c','d','e',0,1,2,3)
print(c[-9])
\end{verbatim}



\begin{center}
\textbf{Завдання 8}
\end{center}

\begin{verbatim}
c = ('0','1','2','3','4','5','6','7','8','9')
print(c[:6:3])
\end{verbatim}



\begin{center}
\textbf{Завдання 9}
\end{center}

\begin{verbatim}
a=['0','1','2','3']
a.remove(3)
print(a)
\end{verbatim}



\begin{center}
\textbf{Завдання 10}
\end{center}
Що буде надруковано за виконання?\par
\begin{verbatim}
c=['0','1','2','3']
del c[4:5]
print(c)
\end{verbatim}

\end{minipage}

%\end {large}
}
\newpage

{{ 27.11.2025 Контрольна робота \No 2, варіант  \No \textbf { 19 }\par}}
{
%\begin {large}

{%Завдання 1-10: Що буде виведено за виконання (повідомлення інтерпретатора про помилку %позначати error)?

%Завдання 11-12 (на звороті): написати функцію для розв'язання задачі. Оцінюється економність обчислень та якість коду.

%Тематика задач: використання циклів, програмування з використанням рекурентних співвідношень,
%обробка списків та рядків.

У завданнях 1-10 вказати, що буде виведено за виконання.

Повідомлення інтерпретатора про помилку позначати ERROR

Завдання 11 та 12 на звороті.

\begin{minipage}[t]{7.5cm}

\begin{center}
\textbf{Завдання 1}
\end{center}

\begin{verbatim}
a, b, c = 6, 7, 6
def f(a, b=9, c=8):
    print(a, b, c, end=" ")

f(b=2, c=4, 0)
print(a, b, c)
\end{verbatim}



\begin{center}
\textbf{Завдання 2}
\end{center}

\begin{verbatim}
a,b,c=7,2,9
def f(b):
    global c
    a=5
    b=1
    c=4
    return a+b+c

a,b,c=0,8,6
print(f(b),a,b,c)
\end{verbatim}



\begin{center}
\textbf{Завдання 3}
\end{center}

\begin{verbatim}
def f():
    try:
        res = int("d0")
    except ValueError: 
        return 4
    except TypeError: 
        return 6
    else: return 35
    finally: return 22
    return res

print(f())
\end{verbatim}



\begin{center}
\textbf{Завдання 4}
\end{center}

\begin{verbatim}
k=4
while k<=10:
    k+=1
print(k)
\end{verbatim}

\end{minipage}
\vline \hspace{6mm}
\begin{minipage}[t]{10.5cm}

\begin{center}
\textbf{Завдання 5}
\end{center}

\begin{verbatim}
for f in range(-4,-4+3,-3):
    if f>=8:
        break
    print(f, end=' ')
    f=-4
else:
    print(13, end=' ')
print(f, end=' ')
\end{verbatim}



\begin{center}
\textbf{Завдання 6}
\end{center}

\begin{verbatim}
def f(n):
    if n > 2:
        f(n // 2)
    else:
        print(n % 2, end=' ')

f(16)
\end{verbatim}



\begin{center}
\textbf{Завдання 7}
\end{center}

\begin{verbatim}
g = ('a','b','c','d','e',0,1,2)
print(g[14])
\end{verbatim}



\begin{center}
\textbf{Завдання 8}
\end{center}

\begin{verbatim}
e = ['0','1','2','3','4','5','6','7','8','9']
print(e[::-2])
\end{verbatim}



\begin{center}
\textbf{Завдання 9}
\end{center}

\begin{verbatim}
h=['0','1','2','3']
h.append('12')
print(h)
\end{verbatim}



\begin{center}
\textbf{Завдання 10}
\end{center}
Що буде надруковано за виконання?\par
\begin{verbatim}
a=['0','1','2','3']
del a[1:4]
print(a)
\end{verbatim}

\end{minipage}

%\end {large}
}
\newpage

{{ 27.11.2025 Контрольна робота \No 2, варіант  \No \textbf { 20 }\par}}
{
%\begin {large}

{%Завдання 1-10: Що буде виведено за виконання (повідомлення інтерпретатора про помилку %позначати error)?

%Завдання 11-12 (на звороті): написати функцію для розв'язання задачі. Оцінюється економність обчислень та якість коду.

%Тематика задач: використання циклів, програмування з використанням рекурентних співвідношень,
%обробка списків та рядків.

У завданнях 1-10 вказати, що буде виведено за виконання.

Повідомлення інтерпретатора про помилку позначати ERROR

Завдання 11 та 12 на звороті.

\begin{minipage}[t]{7.5cm}

\begin{center}
\textbf{Завдання 1}
\end{center}

\begin{verbatim}
a, b, c = 7, 6, 8
def g(a, b=9, c):
    print(a, b, c, end=" ")

g(1, 5, 2)
print(a, b, c)
\end{verbatim}



\begin{center}
\textbf{Завдання 2}
\end{center}

\begin{verbatim}
a,b,c=0,2,8
def h(a):
    global c
    a=4
    b-=1
    c=5
    return a+b+c

a,b,c=9,3,4
print(h(a),a,b,c)
\end{verbatim}



\begin{center}
\textbf{Завдання 3}
\end{center}

\begin{verbatim}
def f():
    try:
        res = int(9%1)
        return 42
    except ZeroDivisionError: 
        return 1
    except BaseException: 
        return 2
    return res

print(f())
\end{verbatim}



\begin{center}
\textbf{Завдання 4}
\end{center}

\begin{verbatim}
m=6
while m>1:
    m+=-2
print(m)
\end{verbatim}

\end{minipage}
\vline \hspace{6mm}
\begin{minipage}[t]{10.5cm}

\begin{center}
\textbf{Завдання 5}
\end{center}

\begin{verbatim}
for b in range(7,7+5,-2):
    if b<=6:
        continue
    print(b, end=' ')
    b=10
else:
    print(b, end=' ')
print(b, end=' ')
\end{verbatim}



\begin{center}
\textbf{Завдання 6}
\end{center}

\begin{verbatim}
def f(s):
    for i in range(len(s)):
        s[i]+=1
    return

s=[1,2,3]
f(s)
print(s)
\end{verbatim}



\begin{center}
\textbf{Завдання 7}
\end{center}

\begin{verbatim}
h = ('a','b','c','d','e',0,1,2,3,4)
print(h[1])
\end{verbatim}



\begin{center}
\textbf{Завдання 8}
\end{center}

\begin{verbatim}
d = '0123456789'
print(d[12:-4:-1])
\end{verbatim}



\begin{center}
\textbf{Завдання 9}
\end{center}

\begin{verbatim}
f=[10,11,12,13]
print(f.index(9), end=' ')
\end{verbatim}



\begin{center}
\textbf{Завдання 10}
\end{center}
Що буде надруковано за виконання?\par
\begin{verbatim}
d=[10,11,12,13]
d[2:-1]='12'
print(d)
\end{verbatim}

\end{minipage}

%\end {large}
}
\newpage

{{ 27.11.2025 Контрольна робота \No 2, варіант  \No \textbf { 21 }\par}}
{
%\begin {large}

{%Завдання 1-10: Що буде виведено за виконання (повідомлення інтерпретатора про помилку %позначати error)?

%Завдання 11-12 (на звороті): написати функцію для розв'язання задачі. Оцінюється економність обчислень та якість коду.

%Тематика задач: використання циклів, програмування з використанням рекурентних співвідношень,
%обробка списків та рядків.

У завданнях 1-10 вказати, що буде виведено за виконання.

Повідомлення інтерпретатора про помилку позначати ERROR

Завдання 11 та 12 на звороті.

\begin{minipage}[t]{7.5cm}

\begin{center}
\textbf{Завдання 1}
\end{center}

\begin{verbatim}
a, b, c = 6, 9, 8
def h(a, b, c=7):
    print(a, b, c, end=" ")

h(c=0, b=1, a=3)
print(a, b, c)
\end{verbatim}



\begin{center}
\textbf{Завдання 2}
\end{center}

\begin{verbatim}
a,b,c=8,5,0
def g(a):
    a=1
    b=3
    c=2
    return a+b+c

a,b,c=3,4,7
print(g(b),a,b,c)
\end{verbatim}



\begin{center}
\textbf{Завдання 3}
\end{center}

\begin{verbatim}
def f():
    try:
        res = int("9")
    except Exception: 
        return 0
    except KeyboardInterrupt: 
        return 4
    finally: return 20
    return res

print(f())
\end{verbatim}



\begin{center}
\textbf{Завдання 4}
\end{center}

\begin{verbatim}
m=6
while m<8:
    m+=2
print(m)
\end{verbatim}

\end{minipage}
\vline \hspace{6mm}
\begin{minipage}[t]{10.5cm}

\begin{center}
\textbf{Завдання 5}
\end{center}

\begin{verbatim}
for e in range(-1,-1+2,2):
    if e<3:
        break
    print(e, end=' ')
    e=-10
    if e>3:
        break
else:
    print(e, end=' ')
print(e, end=' ')
\end{verbatim}



\begin{center}
\textbf{Завдання 6}
\end{center}

\begin{verbatim}
def f(n):
    if n > 9:
        return f(n // 100) + 1
    else:
        return 1

print(f(123456))
\end{verbatim}



\begin{center}
\textbf{Завдання 7}
\end{center}

\begin{verbatim}
b = ('a','b','c','d','e',0,1,2)
print(b[-4])
\end{verbatim}



\begin{center}
\textbf{Завдання 8}
\end{center}

\begin{verbatim}
g = (0,1,2,3,4,5,6,7,8,9)
print(g[::2])
\end{verbatim}



\begin{center}
\textbf{Завдання 9}
\end{center}

\begin{verbatim}
g=[10,11,12,13]
print(g.pop(1), end=' ')
print(g)
\end{verbatim}



\begin{center}
\textbf{Завдання 10}
\end{center}
Що буде надруковано за виконання?\par
\begin{verbatim}
f=['0','1','2','3']
f[-1:]=[4,5]
print(f)
\end{verbatim}

\end{minipage}

%\end {large}
}
\newpage

{{ 27.11.2025 Контрольна робота \No 2, варіант  \No \textbf { 22 }\par}}
{
%\begin {large}

{%Завдання 1-10: Що буде виведено за виконання (повідомлення інтерпретатора про помилку %позначати error)?

%Завдання 11-12 (на звороті): написати функцію для розв'язання задачі. Оцінюється економність обчислень та якість коду.

%Тематика задач: використання циклів, програмування з використанням рекурентних співвідношень,
%обробка списків та рядків.

У завданнях 1-10 вказати, що буде виведено за виконання.

Повідомлення інтерпретатора про помилку позначати ERROR

Завдання 11 та 12 на звороті.

\begin{minipage}[t]{7.5cm}

\begin{center}
\textbf{Завдання 1}
\end{center}

\begin{verbatim}
a, b, c = 6, 7, 8
def f(a, b, c):
    print(a, b, c, end=" ")

f(4, 5, c=2)
print(a, b, c)
\end{verbatim}



\begin{center}
\textbf{Завдання 2}
\end{center}

\begin{verbatim}
a,b,c=2,6,4
def f(b):
    a+=3
    b=4
    c=2
    return a+b+c

a,b,c=8,3,0
print(f(b),a,b,c)
\end{verbatim}



\begin{center}
\textbf{Завдання 3}
\end{center}

\begin{verbatim}
def f():
    try:
        res = int(6//0.0)
        return 41
    except ValueError: 
        return 1
    except BaseException: 
        return 5
    else: return 33
    return res

print(f())
\end{verbatim}



\begin{center}
\textbf{Завдання 4}
\end{center}

\begin{verbatim}
j=7
while j<11:
    j+=2
print(j)
\end{verbatim}

\end{minipage}
\vline \hspace{6mm}
\begin{minipage}[t]{10.5cm}

\begin{center}
\textbf{Завдання 5}
\end{center}

\begin{verbatim}
for a in range(-4,-4+4,-3):
    if a>4:
        continue
    print(a, end=' ')
    a=4
else:
    print(a, end=' ')
print(a, end=' ')
\end{verbatim}



\begin{center}
\textbf{Завдання 6}
\end{center}

\begin{verbatim}
def f(s1):
    s=s1[:]
    for i in range(len(s)):
        s[i]+=1

s=[1,2,3]
f(s)
print(s)
\end{verbatim}



\begin{center}
\textbf{Завдання 7}
\end{center}

\begin{verbatim}
a = ('a','b','c','d','e',0,1,2,3)
print(a[-14])
\end{verbatim}



\begin{center}
\textbf{Завдання 8}
\end{center}

\begin{verbatim}
f = [0,1,2,3,4,5,6,7,8,9]
print(f[-14:3:-1])
\end{verbatim}



\begin{center}
\textbf{Завдання 9}
\end{center}

\begin{verbatim}
b=['0','1','2','3']
b.insert(-4, [2,3])
print(b)
\end{verbatim}



\begin{center}
\textbf{Завдання 10}
\end{center}
Що буде надруковано за виконання?\par
\begin{verbatim}
e=[10,11,12,13]
e[6:]=[12]
print(e)
\end{verbatim}

\end{minipage}

%\end {large}
}
\newpage

{{ 27.11.2025 Контрольна робота \No 2, варіант  \No \textbf { 23 }\par}}
{
%\begin {large}

{%Завдання 1-10: Що буде виведено за виконання (повідомлення інтерпретатора про помилку %позначати error)?

%Завдання 11-12 (на звороті): написати функцію для розв'язання задачі. Оцінюється економність обчислень та якість коду.

%Тематика задач: використання циклів, програмування з використанням рекурентних співвідношень,
%обробка списків та рядків.

У завданнях 1-10 вказати, що буде виведено за виконання.

Повідомлення інтерпретатора про помилку позначати ERROR

Завдання 11 та 12 на звороті.

\begin{minipage}[t]{7.5cm}

\begin{center}
\textbf{Завдання 1}
\end{center}

\begin{verbatim}
a, b, c = 9, 6, 9
def h(a, b=7, c=8):
    print(a, b, c, end=" ")

h(3, 4)
print(a, b, c)
\end{verbatim}



\begin{center}
\textbf{Завдання 2}
\end{center}

\begin{verbatim}
a,b,c=9,5,7
def h(a):
    global c
    a*=1
    b=5
    c=3
    return a+b+c

a,b,c=1,9,3
print(h(a),a,b,c)
\end{verbatim}



\begin{center}
\textbf{Завдання 3}
\end{center}

\begin{verbatim}
def f():
    try:
        res = 3!=7
        return 40
    except KeyboardInterrupt: 
        return 7
    except TypeError: 
        return 2
    return res

print(f())
\end{verbatim}



\begin{center}
\textbf{Завдання 4}
\end{center}

\begin{verbatim}
n=8
while n<=7:
    n+=1
print(n)
\end{verbatim}

\end{minipage}
\vline \hspace{6mm}
\begin{minipage}[t]{10.5cm}

\begin{center}
\textbf{Завдання 5}
\end{center}

\begin{verbatim}
for c in range(9,9+6,3):
    if c<=7:
        continue
    print(c, end=' ')
    c=-4
else:
    print(c, end=' ')
print(c, end=' ')
\end{verbatim}



\begin{center}
\textbf{Завдання 6}
\end{center}

\begin{verbatim}
def f(s):
    s1=s
    for i in range(len(s)):
        s1[i]+=1
    return s1

s=[1,2,3]
f(s)
print(s)
\end{verbatim}



\begin{center}
\textbf{Завдання 7}
\end{center}

\begin{verbatim}
f = ('a','b','c','d','e',0,1,2,4)
print(f[-2])
\end{verbatim}



\begin{center}
\textbf{Завдання 8}
\end{center}

\begin{verbatim}
a = '0123456789'
print(a[5::3])
\end{verbatim}



\begin{center}
\textbf{Завдання 9}
\end{center}

\begin{verbatim}
c=['0','1','2','3']
c.append(19)
print(c)
\end{verbatim}



\begin{center}
\textbf{Завдання 10}
\end{center}
Що буде надруковано за виконання?\par
\begin{verbatim}
h=['0','1','2','3']
del h[:0]
print(h)
\end{verbatim}

\end{minipage}

%\end {large}
}
\newpage

{{ 27.11.2025 Контрольна робота \No 2, варіант  \No \textbf { 24 }\par}}
{
%\begin {large}

{%Завдання 1-10: Що буде виведено за виконання (повідомлення інтерпретатора про помилку %позначати error)?

%Завдання 11-12 (на звороті): написати функцію для розв'язання задачі. Оцінюється економність обчислень та якість коду.

%Тематика задач: використання циклів, програмування з використанням рекурентних співвідношень,
%обробка списків та рядків.

У завданнях 1-10 вказати, що буде виведено за виконання.

Повідомлення інтерпретатора про помилку позначати ERROR

Завдання 11 та 12 на звороті.

\begin{minipage}[t]{7.5cm}

\begin{center}
\textbf{Завдання 1}
\end{center}

\begin{verbatim}
a, b, c = 7, 9, 6
def g(a, b=8, c):
    print(a, b, c, end=" ")

g(a=5, 1, b=0)
print(a, b, c)
\end{verbatim}



\begin{center}
\textbf{Завдання 2}
\end{center}

\begin{verbatim}
a,b,c=9,6,2
def f(b):
    a=4
    b=5
    c=1
    return a+b+c

a,b,c=1,9,1
print(f(a),a,b,c)
\end{verbatim}



\begin{center}
\textbf{Завдання 3}
\end{center}

\begin{verbatim}
def f():
    try:
        res = int(8/1)
    except Exception: 
        return 3
    except ZeroDivisionError: 
        return 9
    else: return 34
    finally: return 25
    return res

print(f())
\end{verbatim}



\begin{center}
\textbf{Завдання 4}
\end{center}

\begin{verbatim}
m=1
while m<=9:
    m+=2
print(m)
\end{verbatim}

\end{minipage}
\vline \hspace{6mm}
\begin{minipage}[t]{10.5cm}

\begin{center}
\textbf{Завдання 5}
\end{center}

\begin{verbatim}
for d in range(-6,-6+5,-2):
    if d>=5:
        pass
    print(d, end=' ')
    d=8
else:
    print(d, end=' ')
print(d, end=' ')
\end{verbatim}



\begin{center}
\textbf{Завдання 6}
\end{center}

\begin{verbatim}
def f(s1):
    s=s1[:]
    for i in range(len(s)):
        s[i]+=1

s=[1,2,3]
f(s)
print(s)
\end{verbatim}



\begin{center}
\textbf{Завдання 7}
\end{center}

\begin{verbatim}
e = ('a','b','c','d','e',0,1)
print(e[4])
\end{verbatim}



\begin{center}
\textbf{Завдання 8}
\end{center}

\begin{verbatim}
h = ['0','1','2','3','4','5','6','7','8','9']
print(h[:13:-3])
\end{verbatim}



\begin{center}
\textbf{Завдання 9}
\end{center}

\begin{verbatim}
d=[10,11,12,13]
d.extend('10')
print(d)
\end{verbatim}



\begin{center}
\textbf{Завдання 10}
\end{center}
Що буде надруковано за виконання?\par
\begin{verbatim}
g=[10,11,12,13]
g[-1:7]=[2,3]
print(g)
\end{verbatim}

\end{minipage}

%\end {large}
}
\newpage

{{ 27.11.2025 Контрольна робота \No 2, варіант  \No \textbf { 25 }\par}}
{
%\begin {large}

{%Завдання 1-10: Що буде виведено за виконання (повідомлення інтерпретатора про помилку %позначати error)?

%Завдання 11-12 (на звороті): написати функцію для розв'язання задачі. Оцінюється економність обчислень та якість коду.

%Тематика задач: використання циклів, програмування з використанням рекурентних співвідношень,
%обробка списків та рядків.

У завданнях 1-10 вказати, що буде виведено за виконання.

Повідомлення інтерпретатора про помилку позначати ERROR

Завдання 11 та 12 на звороті.

\begin{minipage}[t]{7.5cm}

\begin{center}
\textbf{Завдання 1}
\end{center}

\begin{verbatim}
a, b, c = 6, 9, 8
def h(a, b=7, c):
    print(a, b, c, end=" ")

h(1, 0)
print(a, b, c)
\end{verbatim}



\begin{center}
\textbf{Завдання 2}
\end{center}

\begin{verbatim}
a,b,c=8,6,0
def g(b):
    a=4
    b-=2
    c=1
    return a+b+c

a,b,c=5,4,7
print(g(b),a,b,c)
\end{verbatim}



\begin{center}
\textbf{Завдання 3}
\end{center}

\begin{verbatim}
def f():
    try:
        res = int("c4")
        return 45
    except ZeroDivisionError: 
        return 2
    except Exception: 
        return 5
    else: return 30
    finally: return 21
    return res

print(f())
\end{verbatim}



\begin{center}
\textbf{Завдання 4}
\end{center}

\begin{verbatim}
n=0
while n<12:
    n+=1
print(n)
\end{verbatim}

\end{minipage}
\vline \hspace{6mm}
\begin{minipage}[t]{10.5cm}

\begin{center}
\textbf{Завдання 5}
\end{center}

\begin{verbatim}
for d in range(-7,-7+3,3):
    if d<5:
        continue
    print(d, end=' ')
    d=6
else:
    print(d, end=' ')
print(d, end=' ')
\end{verbatim}



\begin{center}
\textbf{Завдання 6}
\end{center}

\begin{verbatim}
def f(s):
    for i in range(len(s)):
        s[i]+=1
    return

s=[1,2,3]
f(s)
print(s)
\end{verbatim}



\begin{center}
\textbf{Завдання 7}
\end{center}

\begin{verbatim}
b = ('a','b','c','d','e',0,1,2,3,4)
print(b[7])
\end{verbatim}



\begin{center}
\textbf{Завдання 8}
\end{center}

\begin{verbatim}
h = ('0','1','2','3','4','5','6','7','8','9')
print(h[:-9:-2])
\end{verbatim}



\begin{center}
\textbf{Завдання 9}
\end{center}

\begin{verbatim}
h=[10,11,12,13]
print(h.index('11'), end=' ')
\end{verbatim}



\begin{center}
\textbf{Завдання 10}
\end{center}
Що буде надруковано за виконання?\par
\begin{verbatim}
b=[10,11,12,13]
b[-6:-4]='11'
print(b)
\end{verbatim}

\end{minipage}

%\end {large}
}
\newpage

{{ 27.11.2025 Контрольна робота \No 2, варіант  \No \textbf { 26 }\par}}
{
%\begin {large}

{%Завдання 1-10: Що буде виведено за виконання (повідомлення інтерпретатора про помилку %позначати error)?

%Завдання 11-12 (на звороті): написати функцію для розв'язання задачі. Оцінюється економність обчислень та якість коду.

%Тематика задач: використання циклів, програмування з використанням рекурентних співвідношень,
%обробка списків та рядків.

У завданнях 1-10 вказати, що буде виведено за виконання.

Повідомлення інтерпретатора про помилку позначати ERROR

Завдання 11 та 12 на звороті.

\begin{minipage}[t]{7.5cm}

\begin{center}
\textbf{Завдання 1}
\end{center}

\begin{verbatim}
a, b, c = 8, 9, 7
def g(a, b, c=6):
    print(a, b, c, end=" ")

g(a=3, 5, c=2)
print(a, b, c)
\end{verbatim}



\begin{center}
\textbf{Завдання 2}
\end{center}

\begin{verbatim}
a,b,c=1,2,4
def f(b):
    a=5
    b+=5
    c=4
    return a+b+c

a,b,c=6,9,1
print(f(a),a,b,c)
\end{verbatim}



\begin{center}
\textbf{Завдання 3}
\end{center}

\begin{verbatim}
def f():
    try:
        res = 0==6
    except ValueError: 
        return 6
    except TypeError: 
        return 7
    return res

print(f())
\end{verbatim}



\begin{center}
\textbf{Завдання 4}
\end{center}

\begin{verbatim}
i=2
while i<5:
    i+=2
print(i)
\end{verbatim}

\end{minipage}
\vline \hspace{6mm}
\begin{minipage}[t]{10.5cm}

\begin{center}
\textbf{Завдання 5}
\end{center}

\begin{verbatim}
for b in range(-5,-5+2,-2):
    if b<=8:
        break
    print(b, end=' ')
    b=9
else:
    print(b, end=' ')
print(b, end=' ')
\end{verbatim}



\begin{center}
\textbf{Завдання 6}
\end{center}

\begin{verbatim}
def f(s):
    for el in s:
        el+=1
    return s

s=[1,2,3]
s=f(s)
print(s)
\end{verbatim}



\begin{center}
\textbf{Завдання 7}
\end{center}

\begin{verbatim}
f = ('a','b','c','d','e',0,1)
print(f[-7])
\end{verbatim}



\begin{center}
\textbf{Завдання 8}
\end{center}

\begin{verbatim}
g = (0,1,2,3,4,5,6,7,8,9)
print(g[6:-2:-2])
\end{verbatim}



\begin{center}
\textbf{Завдання 9}
\end{center}

\begin{verbatim}
g=['0','1','2','3']
g+=[1,0]
print(g)
\end{verbatim}



\begin{center}
\textbf{Завдання 10}
\end{center}
Що буде надруковано за виконання?\par
\begin{verbatim}
d=['0','1','2','3']
del d[-5:-8]
print(d)
\end{verbatim}

\end{minipage}

%\end {large}
}
\newpage

{{ 27.11.2025 Контрольна робота \No 2, варіант  \No \textbf { 27 }\par}}
{
%\begin {large}

{%Завдання 1-10: Що буде виведено за виконання (повідомлення інтерпретатора про помилку %позначати error)?

%Завдання 11-12 (на звороті): написати функцію для розв'язання задачі. Оцінюється економність обчислень та якість коду.

%Тематика задач: використання циклів, програмування з використанням рекурентних співвідношень,
%обробка списків та рядків.

У завданнях 1-10 вказати, що буде виведено за виконання.

Повідомлення інтерпретатора про помилку позначати ERROR

Завдання 11 та 12 на звороті.

\begin{minipage}[t]{7.5cm}

\begin{center}
\textbf{Завдання 1}
\end{center}

\begin{verbatim}
a, b, c = 9, 6, 7
def f(a, b, c):
    print(a, b, c, end=" ")

f(4, 2, 3)
print(a, b, c)
\end{verbatim}



\begin{center}
\textbf{Завдання 2}
\end{center}

\begin{verbatim}
a,b,c=8,5,0
def h(a):
    global c
    a=5
    b=2
    c=4
    return a+b+c

a,b,c=6,3,7
print(h(a),a,b,c)
\end{verbatim}



\begin{center}
\textbf{Завдання 3}
\end{center}

\begin{verbatim}
def f():
    try:
        res = int(1//2)
    except KeyboardInterrupt: 
        return 8
    except BaseException: 
        return 7
    else: return 32
    return res

print(f())
\end{verbatim}



\begin{center}
\textbf{Завдання 4}
\end{center}

\begin{verbatim}
k=9
while k<=13:
    k+=1
print(k)
\end{verbatim}

\end{minipage}
\vline \hspace{6mm}
\begin{minipage}[t]{10.5cm}

\begin{center}
\textbf{Завдання 5}
\end{center}

\begin{verbatim}
for c in range(1,1+6,-3):
    if c>7:
        continue
    print(c, end=' ')
    c=-10
else:
    print(c, end=' ')
print(c, end=' ')
\end{verbatim}



\begin{center}
\textbf{Завдання 6}
\end{center}

\begin{verbatim}
def f(n):
    if n > 9:
        f(n // 10)
    else:
        print(n % 10, end=' ')

f(543216)
\end{verbatim}



\begin{center}
\textbf{Завдання 7}
\end{center}

\begin{verbatim}
g = ('a','b','c','d','e',0,1,2,4)
print(g[2])
\end{verbatim}



\begin{center}
\textbf{Завдання 8}
\end{center}

\begin{verbatim}
a = [0,1,2,3,4,5,6,7,8,9]
print(a[1:-15:2])
\end{verbatim}



\begin{center}
\textbf{Завдання 9}
\end{center}

\begin{verbatim}
a=[10,11,12,13]
print(a.pop(5), end=' ')
print(a)
\end{verbatim}



\begin{center}
\textbf{Завдання 10}
\end{center}
Що буде надруковано за виконання?\par
\begin{verbatim}
e=['0','1','2','3']
e[4:6]=(15)
print(e)
\end{verbatim}

\end{minipage}

%\end {large}
}
\newpage

{{ 27.11.2025 Контрольна робота \No 2, варіант  \No \textbf { 28 }\par}}
{
%\begin {large}

{%Завдання 1-10: Що буде виведено за виконання (повідомлення інтерпретатора про помилку %позначати error)?

%Завдання 11-12 (на звороті): написати функцію для розв'язання задачі. Оцінюється економність обчислень та якість коду.

%Тематика задач: використання циклів, програмування з використанням рекурентних співвідношень,
%обробка списків та рядків.

У завданнях 1-10 вказати, що буде виведено за виконання.

Повідомлення інтерпретатора про помилку позначати ERROR

Завдання 11 та 12 на звороті.

\begin{minipage}[t]{7.5cm}

\begin{center}
\textbf{Завдання 1}
\end{center}

\begin{verbatim}
a, b, c = 8, 9, 6
def f(a, b=8, c=7):
    print(a, b, c, end=" ")

f(b=0, c=4, 1)
print(a, b, c)
\end{verbatim}



\begin{center}
\textbf{Завдання 2}
\end{center}

\begin{verbatim}
a,b,c=3,5,7
def h(a):
    global c
    a-=1
    b=3
    c=2
    return a+b+c

a,b,c=2,0,8
print(h(a),a,b,c)
\end{verbatim}



\begin{center}
\textbf{Завдання 3}
\end{center}

\begin{verbatim}
def f():
    try:
        res = int("8")
        return 41
    except Exception: 
        return 3
    except TypeError: 
        return 5
    finally: return 23
    return res

print(f())
\end{verbatim}



\begin{center}
\textbf{Завдання 4}
\end{center}

\begin{verbatim}
n=8
while n>=5:
    n+=-3
print(n)
\end{verbatim}

\end{minipage}
\vline \hspace{6mm}
\begin{minipage}[t]{10.5cm}

\begin{center}
\textbf{Завдання 5}
\end{center}

\begin{verbatim}
for a in range(0,0+5,2):
    if a>=6:
        pass
    print(a, end=' ')
    a=3
else:
    print(13, end=' ')
print(a, end=' ')
\end{verbatim}



\begin{center}
\textbf{Завдання 6}
\end{center}

\begin{verbatim}
def f(s):
    s1=s
    for i in range(len(s)):
        s1[i]+=1
    return s1

s=[1,2,3]
f(s)
print(s)
\end{verbatim}



\begin{center}
\textbf{Завдання 7}
\end{center}

\begin{verbatim}
d = ('a','b','c','d','e',0,1,2,3)
print(d[-3])
\end{verbatim}



\begin{center}
\textbf{Завдання 8}
\end{center}

\begin{verbatim}
b = ('0','1','2','3','4','5','6','7','8','9')
print(b[-13:-7:-3])
\end{verbatim}



\begin{center}
\textbf{Завдання 9}
\end{center}

\begin{verbatim}
e=['0','1','2','3']
e.remove(1)
print(e)
\end{verbatim}



\begin{center}
\textbf{Завдання 10}
\end{center}
Що буде надруковано за виконання?\par
\begin{verbatim}
a=[10,11,12,13]
del a[-2]
print(a)
\end{verbatim}

\end{minipage}

%\end {large}
}
\newpage

{{ 27.11.2025 Контрольна робота \No 2, варіант  \No \textbf { 29 }\par}}
{
%\begin {large}

{%Завдання 1-10: Що буде виведено за виконання (повідомлення інтерпретатора про помилку %позначати error)?

%Завдання 11-12 (на звороті): написати функцію для розв'язання задачі. Оцінюється економність обчислень та якість коду.

%Тематика задач: використання циклів, програмування з використанням рекурентних співвідношень,
%обробка списків та рядків.

У завданнях 1-10 вказати, що буде виведено за виконання.

Повідомлення інтерпретатора про помилку позначати ERROR

Завдання 11 та 12 на звороті.

\begin{minipage}[t]{7.5cm}

\begin{center}
\textbf{Завдання 1}
\end{center}

\begin{verbatim}
a, b, c = 7, 8, 9
def g(a, b, c):
    print(a, b, c, end=" ")

g(b=5, a=2, c=5)
print(a, b, c)
\end{verbatim}



\begin{center}
\textbf{Завдання 2}
\end{center}

\begin{verbatim}
a,b,c=4,2,7
def g(b):
    a=3
    b=2
    c=5
    return a+b+c

a,b,c=8,5,2
print(g(b),a,b,c)
\end{verbatim}



\begin{center}
\textbf{Завдання 3}
\end{center}

\begin{verbatim}
def f():
    try:
        res = int("d6")
    except ValueError: 
        return 0
    except BaseException: 
        return 1
    else: return 31
    finally: return 22
    return res

print(f())
\end{verbatim}



\begin{center}
\textbf{Завдання 4}
\end{center}

\begin{verbatim}
t=9
while t>9:
    t+=-3
print(t)
\end{verbatim}

\end{minipage}
\vline \hspace{6mm}
\begin{minipage}[t]{10.5cm}

\begin{center}
\textbf{Завдання 5}
\end{center}

\begin{verbatim}
for f in range(2,2+3,3):
    if f>3:
        continue
    print(f, end=' ')
    f=-1
else:
    print(f, end=' ')
print(f, end=' ')
\end{verbatim}



\begin{center}
\textbf{Завдання 6}
\end{center}

\begin{verbatim}
def f(n):
    if n > 2:
        f(n // 2)
    print(n % 2, end=' ')
    

f(16)
\end{verbatim}



\begin{center}
\textbf{Завдання 7}
\end{center}

\begin{verbatim}
e = ('a','b','c','d','e',0,1,2)
print(e[-6])
\end{verbatim}



\begin{center}
\textbf{Завдання 8}
\end{center}

\begin{verbatim}
c = '0123456789'
print(c[-9::-1])
\end{verbatim}



\begin{center}
\textbf{Завдання 9}
\end{center}

\begin{verbatim}
c=['0','1','2','3']
c.append('13')
print(c)
\end{verbatim}



\begin{center}
\textbf{Завдання 10}
\end{center}
Що буде надруковано за виконання?\par
\begin{verbatim}
g=[10,11,12,13]
g[:4]='10'
print(g)
\end{verbatim}

\end{minipage}

%\end {large}
}
\newpage

{{ 27.11.2025 Контрольна робота \No 2, варіант  \No \textbf { 30 }\par}}
{
%\begin {large}

{%Завдання 1-10: Що буде виведено за виконання (повідомлення інтерпретатора про помилку %позначати error)?

%Завдання 11-12 (на звороті): написати функцію для розв'язання задачі. Оцінюється економність обчислень та якість коду.

%Тематика задач: використання циклів, програмування з використанням рекурентних співвідношень,
%обробка списків та рядків.

У завданнях 1-10 вказати, що буде виведено за виконання.

Повідомлення інтерпретатора про помилку позначати ERROR

Завдання 11 та 12 на звороті.

\begin{minipage}[t]{7.5cm}

\begin{center}
\textbf{Завдання 1}
\end{center}

\begin{verbatim}
a, b, c = 6, 7, 8
def h(a, b=9, c):
    print(a, b, c, end=" ")

h(1, a=3)
print(a, b, c)
\end{verbatim}



\begin{center}
\textbf{Завдання 2}
\end{center}

\begin{verbatim}
a,b,c=0,3,5
def g(b):
    a*=4
    b=1
    c=5
    return a+b+c

a,b,c=6,2,1
print(g(b),a,b,c)
\end{verbatim}



\begin{center}
\textbf{Завдання 3}
\end{center}

\begin{verbatim}
def f():
    try:
        res = int(4%0)
        return 44
    except ZeroDivisionError: 
        return 2
    except KeyboardInterrupt: 
        return 9
    return res

print(f())
\end{verbatim}



\begin{center}
\textbf{Завдання 4}
\end{center}

\begin{verbatim}
j=4
while j<6:
    j+=2
print(j)
\end{verbatim}

\end{minipage}
\vline \hspace{6mm}
\begin{minipage}[t]{10.5cm}

\begin{center}
\textbf{Завдання 5}
\end{center}

\begin{verbatim}
for e in range(-1,-1+4,-2):
    if e<=4:
        pass
    print(e, end=' ')
    e=1
    if e<4:
        break
else:
    print(e, end=' ')
print(e, end=' ')
\end{verbatim}



\begin{center}
\textbf{Завдання 6}
\end{center}

\begin{verbatim}
def f(s):
    for i in range(len(s)):
        s[i]+=1
    return

s=[1,2,3]
f(s)
print(s)
\end{verbatim}



\begin{center}
\textbf{Завдання 7}
\end{center}

\begin{verbatim}
a = ('a','b','c','d','e',0,1,2,3,4)
print(a[-11])
\end{verbatim}



\begin{center}
\textbf{Завдання 8}
\end{center}

\begin{verbatim}
d = ['0','1','2','3','4','5','6','7','8','9']
print(d[2:5:3])
\end{verbatim}



\begin{center}
\textbf{Завдання 9}
\end{center}

\begin{verbatim}
f=[10,11,12,13]
print(f.index(10), end=' ')
\end{verbatim}



\begin{center}
\textbf{Завдання 10}
\end{center}
Що буде надруковано за виконання?\par
\begin{verbatim}
c=['0','1','2','3']
c[3:1]='13'
print(c)
\end{verbatim}

\end{minipage}

%\end {large}
}
\newpage

{{ 27.11.2025 Контрольна робота \No 2, варіант  \No \textbf { 31 }\par}}
{
%\begin {large}

{%Завдання 1-10: Що буде виведено за виконання (повідомлення інтерпретатора про помилку %позначати error)?

%Завдання 11-12 (на звороті): написати функцію для розв'язання задачі. Оцінюється економність обчислень та якість коду.

%Тематика задач: використання циклів, програмування з використанням рекурентних співвідношень,
%обробка списків та рядків.

У завданнях 1-10 вказати, що буде виведено за виконання.

Повідомлення інтерпретатора про помилку позначати ERROR

Завдання 11 та 12 на звороті.

\begin{minipage}[t]{7.5cm}

\begin{center}
\textbf{Завдання 1}
\end{center}

\begin{verbatim}
a, b, c = 6, 8, 9
def f(a, b, c=6):
    print(a, b, c, end=" ")

f(4, c=0, b=2)
print(a, b, c)
\end{verbatim}



\begin{center}
\textbf{Завдання 2}
\end{center}

\begin{verbatim}
a,b,c=4,8,9
def h(a):
    global c
    a+=2
    b=3
    c=3
    return a+b+c

a,b,c=7,4,7
print(h(a),a,b,c)
\end{verbatim}



\begin{center}
\textbf{Завдання 3}
\end{center}

\begin{verbatim}
def f():
    try:
        res = int("b7")
        return 40
    except KeyboardInterrupt: 
        return 4
    except ValueError: 
        return 1
    else: return 33
    return res

print(f())
\end{verbatim}



\begin{center}
\textbf{Завдання 4}
\end{center}

\begin{verbatim}
t=5
while t<=10:
    t+=1
print(t)
\end{verbatim}

\end{minipage}
\vline \hspace{6mm}
\begin{minipage}[t]{10.5cm}

\begin{center}
\textbf{Завдання 5}
\end{center}

\begin{verbatim}
for d in range(3,3+2,2):
    if d>=5:
        break
    print(d, end=' ')
    d=-7
else:
    print(d, end=' ')
print(d, end=' ')
\end{verbatim}



\begin{center}
\textbf{Завдання 6}
\end{center}

\begin{verbatim}
def f(s):
    s1=s
    for i in range(len(s)):
        s1[i]+=1
    return s1

s=[1,2,3]
f(s)
print(s)
\end{verbatim}



\begin{center}
\textbf{Завдання 7}
\end{center}

\begin{verbatim}
h = ('a','b','c','d','e',0,1,2,3)
print(h[1])
\end{verbatim}



\begin{center}
\textbf{Завдання 8}
\end{center}

\begin{verbatim}
e = [0,1,2,3,4,5,6,7,8,9]
print(e[:-13:-2])
\end{verbatim}



\begin{center}
\textbf{Завдання 9}
\end{center}

\begin{verbatim}
d=['0','1','2','3']
d+='11'
print(d)
\end{verbatim}



\begin{center}
\textbf{Завдання 10}
\end{center}
Що буде надруковано за виконання?\par
\begin{verbatim}
f=['0','1','2','3']
f[-3:]=()
print(f)
\end{verbatim}

\end{minipage}

%\end {large}
}
\newpage

{{ 27.11.2025 Контрольна робота \No 2, варіант  \No \textbf { 32 }\par}}
{
%\begin {large}

{%Завдання 1-10: Що буде виведено за виконання (повідомлення інтерпретатора про помилку %позначати error)?

%Завдання 11-12 (на звороті): написати функцію для розв'язання задачі. Оцінюється економність обчислень та якість коду.

%Тематика задач: використання циклів, програмування з використанням рекурентних співвідношень,
%обробка списків та рядків.

У завданнях 1-10 вказати, що буде виведено за виконання.

Повідомлення інтерпретатора про помилку позначати ERROR

Завдання 11 та 12 на звороті.

\begin{minipage}[t]{7.5cm}

\begin{center}
\textbf{Завдання 1}
\end{center}

\begin{verbatim}
a, b, c = 7, 7, 9
def h(a, b=8, c=6):
    print(a, b, c, end=" ")

h(4, 5, b=0)
print(a, b, c)
\end{verbatim}



\begin{center}
\textbf{Завдання 2}
\end{center}

\begin{verbatim}
a,b,c=9,8,2
def g(b):
    a=1
    b*=4
    c=5
    return a+b+c

a,b,c=1,6,3
print(g(b),a,b,c)
\end{verbatim}



\begin{center}
\textbf{Завдання 3}
\end{center}

\begin{verbatim}
def f():
    try:
        res = int(8//0)
    except ZeroDivisionError: 
        return 2
    except Exception: 
        return 0
    finally: return 24
    return res

print(f())
\end{verbatim}



\begin{center}
\textbf{Завдання 4}
\end{center}

\begin{verbatim}
k=6
while k<9:
    k+=2
print(k)
\end{verbatim}

\end{minipage}
\vline \hspace{6mm}
\begin{minipage}[t]{10.5cm}

\begin{center}
\textbf{Завдання 5}
\end{center}

\begin{verbatim}
for a in range(-2,-2+6,-3):
    if a<8:
        continue
    print(a, end=' ')
    a=-8
else:
    print(13, end=' ')
print(a, end=' ')
\end{verbatim}



\begin{center}
\textbf{Завдання 6}
\end{center}

\begin{verbatim}
def f(s1):
    s=s1[:]
    for i in range(len(s)):
        s[i]+=1

s=[1,2,3]
f(s)
print(s)
\end{verbatim}



\begin{center}
\textbf{Завдання 7}
\end{center}

\begin{verbatim}
c = ('a','b','c','d','e',0,1)
print(c[-1])
\end{verbatim}



\begin{center}
\textbf{Завдання 8}
\end{center}

\begin{verbatim}
f = '0123456789'
print(f[4:12:2])
\end{verbatim}



\begin{center}
\textbf{Завдання 9}
\end{center}

\begin{verbatim}
b=[10,11,12,13]
b.insert(2, [1,2])
print(b)
\end{verbatim}



\begin{center}
\textbf{Завдання 10}
\end{center}
Що буде надруковано за виконання?\par
\begin{verbatim}
h=[10,11,12,13]
del h[5]
print(h)
\end{verbatim}

\end{minipage}

%\end {large}
}
\newpage

{{ 27.11.2025 Контрольна робота \No 2, варіант  \No \textbf { 33 }\par}}
{
%\begin {large}

{%Завдання 1-10: Що буде виведено за виконання (повідомлення інтерпретатора про помилку %позначати error)?

%Завдання 11-12 (на звороті): написати функцію для розв'язання задачі. Оцінюється економність обчислень та якість коду.

%Тематика задач: використання циклів, програмування з використанням рекурентних співвідношень,
%обробка списків та рядків.

У завданнях 1-10 вказати, що буде виведено за виконання.

Повідомлення інтерпретатора про помилку позначати ERROR

Завдання 11 та 12 на звороті.

\begin{minipage}[t]{7.5cm}

\begin{center}
\textbf{Завдання 1}
\end{center}

\begin{verbatim}
a, b, c = 7, 6, 9
def g(a, b, c):
    print(a, b, c, end=" ")

g(3, c=1)
print(a, b, c)
\end{verbatim}



\begin{center}
\textbf{Завдання 2}
\end{center}

\begin{verbatim}
a,b,c=5,0,4
def f(a):
    global c
    a=2
    b-=2
    c=3
    return a+b+c

a,b,c=1,9,2
print(f(b),a,b,c)
\end{verbatim}



\begin{center}
\textbf{Завдання 3}
\end{center}

\begin{verbatim}
def f():
    try:
        res = int(3/1)
        return 43
    except TypeError: 
        return 6
    except BaseException: 
        return 5
    return res

print(f())
\end{verbatim}



\begin{center}
\textbf{Завдання 4}
\end{center}

\begin{verbatim}
i=3
while i<=11:
    i+=1
print(i)
\end{verbatim}

\end{minipage}
\vline \hspace{6mm}
\begin{minipage}[t]{10.5cm}

\begin{center}
\textbf{Завдання 5}
\end{center}

\begin{verbatim}
for b in range(10,10+3,-2):
    if b>=3:
        continue
    print(b, end=' ')
    b=7
else:
    print(b, end=' ')
print(b, end=' ')
\end{verbatim}



\begin{center}
\textbf{Завдання 6}
\end{center}

\begin{verbatim}
def f(s):
    for el in s:
        el+=1
    return s

s=[1,2,3]
s=f(s)
print(s)
\end{verbatim}



\begin{center}
\textbf{Завдання 7}
\end{center}

\begin{verbatim}
b = ('a','b','c','d','e',0,1,2)
print(b[-7])
\end{verbatim}



\begin{center}
\textbf{Завдання 8}
\end{center}

\begin{verbatim}
h = ('0','1','2','3','4','5','6','7','8','9')
print(h[-2:-8:-3])
\end{verbatim}



\begin{center}
\textbf{Завдання 9}
\end{center}

\begin{verbatim}
d=[10,11,12,13]
d.remove(11)
print(d)
\end{verbatim}



\begin{center}
\textbf{Завдання 10}
\end{center}
Що буде надруковано за виконання?\par
\begin{verbatim}
e=[10,11,12,13]
e[8:-8]=(13,)
print(e)
\end{verbatim}

\end{minipage}

%\end {large}
}
\newpage

{{ 27.11.2025 Контрольна робота \No 2, варіант  \No \textbf { 34 }\par}}
{
%\begin {large}

{%Завдання 1-10: Що буде виведено за виконання (повідомлення інтерпретатора про помилку %позначати error)?

%Завдання 11-12 (на звороті): написати функцію для розв'язання задачі. Оцінюється економність обчислень та якість коду.

%Тематика задач: використання циклів, програмування з використанням рекурентних співвідношень,
%обробка списків та рядків.

У завданнях 1-10 вказати, що буде виведено за виконання.

Повідомлення інтерпретатора про помилку позначати ERROR

Завдання 11 та 12 на звороті.

\begin{minipage}[t]{7.5cm}

\begin{center}
\textbf{Завдання 1}
\end{center}

\begin{verbatim}
a, b, c = 8, 9, 6
def f(a, b=8, c=7):
    print(a, b, c, end=" ")

f(1, 2)
print(a, b, c)
\end{verbatim}



\begin{center}
\textbf{Завдання 2}
\end{center}

\begin{verbatim}
a,b,c=3,5,8
def f(a):
    a-=1
    b=5
    c=4
    return a+b+c

a,b,c=6,0,7
print(f(b),a,b,c)
\end{verbatim}



\begin{center}
\textbf{Завдання 3}
\end{center}

\begin{verbatim}
def f():
    try:
        res = int("9")
    except BaseException: 
        return 8
    except TypeError: 
        return 0
    else: return 35
    finally: return 25
    return res

print(f())
\end{verbatim}



\begin{center}
\textbf{Завдання 4}
\end{center}

\begin{verbatim}
t=7
while t<=13:
    t+=2
print(t)
\end{verbatim}

\end{minipage}
\vline \hspace{6mm}
\begin{minipage}[t]{10.5cm}

\begin{center}
\textbf{Завдання 5}
\end{center}

\begin{verbatim}
for f in range(-10,-10+5,3):
    if f<4:
        pass
    print(f, end=' ')
    f=-5
else:
    print(f, end=' ')
print(f, end=' ')
\end{verbatim}



\begin{center}
\textbf{Завдання 6}
\end{center}

\begin{verbatim}
def f(n):
    if n > 9:
        f(n // 100)
    else:
        print(n % 10, end=' ')

f(123456)
\end{verbatim}



\begin{center}
\textbf{Завдання 7}
\end{center}

\begin{verbatim}
g = ('a','b','c','d','e',0,1,2,4)
print(g[-11])
\end{verbatim}



\begin{center}
\textbf{Завдання 8}
\end{center}

\begin{verbatim}
d = (0,1,2,3,4,5,6,7,8,9)
print(d[:-1:-1])
\end{verbatim}



\begin{center}
\textbf{Завдання 9}
\end{center}

\begin{verbatim}
b=['0','1','2','3']
b.extend(18)
print(b)
\end{verbatim}



\begin{center}
\textbf{Завдання 10}
\end{center}
Що буде надруковано за виконання?\par
\begin{verbatim}
b=['0','1','2','3']
b[:-2]=[1,0]
print(b)
\end{verbatim}

\end{minipage}

%\end {large}
}
\newpage

{{ 27.11.2025 Контрольна робота \No 2, варіант  \No \textbf { 35 }\par}}
{
%\begin {large}

{%Завдання 1-10: Що буде виведено за виконання (повідомлення інтерпретатора про помилку %позначати error)?

%Завдання 11-12 (на звороті): написати функцію для розв'язання задачі. Оцінюється економність обчислень та якість коду.

%Тематика задач: використання циклів, програмування з використанням рекурентних співвідношень,
%обробка списків та рядків.

У завданнях 1-10 вказати, що буде виведено за виконання.

Повідомлення інтерпретатора про помилку позначати ERROR

Завдання 11 та 12 на звороті.

\begin{minipage}[t]{7.5cm}

\begin{center}
\textbf{Завдання 1}
\end{center}

\begin{verbatim}
a, b, c = 8, 7, 9
def g(a, b, c=6):
    print(a, b, c, end=" ")

g(b=3, c=0, 5)
print(a, b, c)
\end{verbatim}



\begin{center}
\textbf{Завдання 2}
\end{center}

\begin{verbatim}
a,b,c=4,7,1
def h(b):
    global c
    a=1
    b*=5
    c=2
    return a+b+c

a,b,c=5,8,9
print(h(a),a,b,c)
\end{verbatim}



\begin{center}
\textbf{Завдання 3}
\end{center}

\begin{verbatim}
def f():
    try:
        res = 7!=2
        return 42
    except ZeroDivisionError: 
        return 4
    except Exception: 
        return 6
    else: return 34
    return res

print(f())
\end{verbatim}



\begin{center}
\textbf{Завдання 4}
\end{center}

\begin{verbatim}
j=8
while j<7:
    j+=1
print(j)
\end{verbatim}

\end{minipage}
\vline \hspace{6mm}
\begin{minipage}[t]{10.5cm}

\begin{center}
\textbf{Завдання 5}
\end{center}

\begin{verbatim}
for c in range(8,8+4,2):
    if c<=6:
        break
    print(c, end=' ')
    c=0
else:
    print(c, end=' ')
print(c, end=' ')
\end{verbatim}



\begin{center}
\textbf{Завдання 6}
\end{center}

\begin{verbatim}
def f(s):
    for i in range(len(s)):
        s[i]+=1
    return

s=[1,2,3]
f(s)
print(s)
\end{verbatim}



\begin{center}
\textbf{Завдання 7}
\end{center}

\begin{verbatim}
f = ('a','b','c','d','e',0,1,2,3,4)
print(f[4])
\end{verbatim}



\begin{center}
\textbf{Завдання 8}
\end{center}

\begin{verbatim}
b = ['0','1','2','3','4','5','6','7','8','9']
print(b[-4::3])
\end{verbatim}



\begin{center}
\textbf{Завдання 9}
\end{center}

\begin{verbatim}
g=['0','1','2','3']
print(g.pop(-1), end=' ')
print(g)
\end{verbatim}



\begin{center}
\textbf{Завдання 10}
\end{center}
Що буде надруковано за виконання?\par
\begin{verbatim}
d=[10,11,12,13]
del d[1:0]
print(d)
\end{verbatim}

\end{minipage}

%\end {large}
}
\newpage

{{ 27.11.2025 Контрольна робота \No 2, варіант  \No \textbf { 36 }\par}}
{
%\begin {large}

{%Завдання 1-10: Що буде виведено за виконання (повідомлення інтерпретатора про помилку %позначати error)?

%Завдання 11-12 (на звороті): написати функцію для розв'язання задачі. Оцінюється економність обчислень та якість коду.

%Тематика задач: використання циклів, програмування з використанням рекурентних співвідношень,
%обробка списків та рядків.

У завданнях 1-10 вказати, що буде виведено за виконання.

Повідомлення інтерпретатора про помилку позначати ERROR

Завдання 11 та 12 на звороті.

\begin{minipage}[t]{7.5cm}

\begin{center}
\textbf{Завдання 1}
\end{center}

\begin{verbatim}
a, b, c = 8, 7, 6
def h(a, b=9, c):
    print(a, b, c, end=" ")

h(b=4, a=2, c=3)
print(a, b, c)
\end{verbatim}



\begin{center}
\textbf{Завдання 2}
\end{center}

\begin{verbatim}
a,b,c=0,3,6
def f(a):
    global c
    a=4
    b=3
    c=1
    return a+b+c

a,b,c=1,9,4
print(f(b),a,b,c)
\end{verbatim}



\begin{center}
\textbf{Завдання 3}
\end{center}

\begin{verbatim}
def f():
    try:
        res = 5<=1
    except ValueError: 
        return 1
    except KeyboardInterrupt: 
        return 2
    finally: return 20
    return res

print(f())
\end{verbatim}



\begin{center}
\textbf{Завдання 4}
\end{center}

\begin{verbatim}
i=10
while i>=4:
    i+=-2
print(i)
\end{verbatim}

\end{minipage}
\vline \hspace{6mm}
\begin{minipage}[t]{10.5cm}

\begin{center}
\textbf{Завдання 5}
\end{center}

\begin{verbatim}
for e in range(-8,-8+2,-3):
    if e>7:
        continue
    print(e, end=' ')
    e=-2
else:
    print(e, end=' ')
print(e, end=' ')
\end{verbatim}



\begin{center}
\textbf{Завдання 6}
\end{center}

\begin{verbatim}
def f(n):
    if n <= 9:
        print(n % 10, end=' ')
    else:
        f(n // 10)

f(987651)
\end{verbatim}



\begin{center}
\textbf{Завдання 7}
\end{center}

\begin{verbatim}
a = ('a','b','c','d','e',0,1,2,4)
print(a[-9])
\end{verbatim}



\begin{center}
\textbf{Завдання 8}
\end{center}

\begin{verbatim}
f = (0,1,2,3,4,5,6,7,8,9)
print(f[-5:-5:2])
\end{verbatim}



\begin{center}
\textbf{Завдання 9}
\end{center}

\begin{verbatim}
a=[10,11,12,13]
a.extend([4,5])
print(a)
\end{verbatim}



\begin{center}
\textbf{Завдання 10}
\end{center}
Що буде надруковано за виконання?\par
\begin{verbatim}
a=['0','1','2','3']
a[2:8]=[14,]
print(a)
\end{verbatim}

\end{minipage}

%\end {large}
}
\newpage

{{ 27.11.2025 Контрольна робота \No 2, варіант  \No \textbf { 37 }\par}}
{
%\begin {large}

{%Завдання 1-10: Що буде виведено за виконання (повідомлення інтерпретатора про помилку %позначати error)?

%Завдання 11-12 (на звороті): написати функцію для розв'язання задачі. Оцінюється економність обчислень та якість коду.

%Тематика задач: використання циклів, програмування з використанням рекурентних співвідношень,
%обробка списків та рядків.

У завданнях 1-10 вказати, що буде виведено за виконання.

Повідомлення інтерпретатора про помилку позначати ERROR

Завдання 11 та 12 на звороті.

\begin{minipage}[t]{7.5cm}

\begin{center}
\textbf{Завдання 1}
\end{center}

\begin{verbatim}
a, b, c = 8, 7, 9
def g(a, b=6, c=8):
    print(a, b, c, end=" ")

g(0, c=4, b=1)
print(a, b, c)
\end{verbatim}



\begin{center}
\textbf{Завдання 2}
\end{center}

\begin{verbatim}
a,b,c=0,6,2
def g(a):
    a+=3
    b=4
    c=5
    return a+b+c

a,b,c=3,2,9
print(g(a),a,b,c)
\end{verbatim}



\begin{center}
\textbf{Завдання 3}
\end{center}

\begin{verbatim}
def f():
    try:
        res = int(3%2)
    except ZeroDivisionError: 
        return 9
    except KeyboardInterrupt: 
        return 0
    else: return 35
    return res

print(f())
\end{verbatim}



\begin{center}
\textbf{Завдання 4}
\end{center}

\begin{verbatim}
n=12
while n>7:
    n+=-2
print(n)
\end{verbatim}

\end{minipage}
\vline \hspace{6mm}
\begin{minipage}[t]{10.5cm}

\begin{center}
\textbf{Завдання 5}
\end{center}

\begin{verbatim}
for a in range(-3,-3+4,-3):
    if a<=4:
        break
    print(a, end=' ')
    a=10
else:
    print(a, end=' ')
print(a, end=' ')
\end{verbatim}



\begin{center}
\textbf{Завдання 6}
\end{center}

\begin{verbatim}
def f(n):
    if n > 9:
        f(n // 10)
    else:
        print(n % 10, end=' ')

f(123456)
\end{verbatim}



\begin{center}
\textbf{Завдання 7}
\end{center}

\begin{verbatim}
c = ('a','b','c','d','e',0,1)
print(c[-4])
\end{verbatim}



\begin{center}
\textbf{Завдання 8}
\end{center}

\begin{verbatim}
c = ['0','1','2','3','4','5','6','7','8','9']
print(c[11:-3:-3])
\end{verbatim}



\begin{center}
\textbf{Завдання 9}
\end{center}

\begin{verbatim}
c=[10,11,12,13]
c.remove('12')
print(c)
\end{verbatim}



\begin{center}
\textbf{Завдання 10}
\end{center}
Що буде надруковано за виконання?\par
\begin{verbatim}
c=[10,11,12,13]
c[0:]=[1,2]
print(c)
\end{verbatim}

\end{minipage}

%\end {large}
}
\newpage

{{ 27.11.2025 Контрольна робота \No 2, варіант  \No \textbf { 38 }\par}}
{
%\begin {large}

{%Завдання 1-10: Що буде виведено за виконання (повідомлення інтерпретатора про помилку %позначати error)?

%Завдання 11-12 (на звороті): написати функцію для розв'язання задачі. Оцінюється економність обчислень та якість коду.

%Тематика задач: використання циклів, програмування з використанням рекурентних співвідношень,
%обробка списків та рядків.

У завданнях 1-10 вказати, що буде виведено за виконання.

Повідомлення інтерпретатора про помилку позначати ERROR

Завдання 11 та 12 на звороті.

\begin{minipage}[t]{7.5cm}

\begin{center}
\textbf{Завдання 1}
\end{center}

\begin{verbatim}
a, b, c = 7, 9, 6
def h(a, b, c):
    print(a, b, c, end=" ")

h(5, 1, 4)
print(a, b, c)
\end{verbatim}



\begin{center}
\textbf{Завдання 2}
\end{center}

\begin{verbatim}
a,b,c=5,6,3
def g(b):
    global c
    a*=4
    b=1
    c=2
    return a+b+c

a,b,c=7,8,0
print(g(b),a,b,c)
\end{verbatim}



\begin{center}
\textbf{Завдання 3}
\end{center}

\begin{verbatim}
def f():
    try:
        res = int("4")
        return 43
    except TypeError: 
        return 9
    except BaseException: 
        return 8
    finally: return 23
    return res

print(f())
\end{verbatim}



\begin{center}
\textbf{Завдання 4}
\end{center}

\begin{verbatim}
j=5
while j>=6:
    j+=-3
print(j)
\end{verbatim}

\end{minipage}
\vline \hspace{6mm}
\begin{minipage}[t]{10.5cm}

\begin{center}
\textbf{Завдання 5}
\end{center}

\begin{verbatim}
for f in range(4,4+6,2):
    if f<8:
        pass
    print(f, end=' ')
    f=5
else:
    print(13, end=' ')
print(f, end=' ')
\end{verbatim}



\begin{center}
\textbf{Завдання 6}
\end{center}

\begin{verbatim}
def f(s):
    for el in s:
        el+=1
    return s

s=[1,2,3]
s=f(s)
print(s)
\end{verbatim}



\begin{center}
\textbf{Завдання 7}
\end{center}

\begin{verbatim}
d = ('a','b','c','d','e',0,1,2,3,4)
print(d[2])
\end{verbatim}



\begin{center}
\textbf{Завдання 8}
\end{center}

\begin{verbatim}
e = [0,1,2,3,4,5,6,7,8,9]
print(e[::-1])
\end{verbatim}



\begin{center}
\textbf{Завдання 9}
\end{center}

\begin{verbatim}
f=['0','1','2','3']
print(f.pop(-5), end=' ')
print(f)
\end{verbatim}



\begin{center}
\textbf{Завдання 10}
\end{center}
Що буде надруковано за виконання?\par
\begin{verbatim}
f=['0','1','2','3']
f[-9:7]='12'
print(f)
\end{verbatim}

\end{minipage}

%\end {large}
}
\newpage

{{ 27.11.2025 Контрольна робота \No 2, варіант  \No \textbf { 39 }\par}}
{
%\begin {large}

{%Завдання 1-10: Що буде виведено за виконання (повідомлення інтерпретатора про помилку %позначати error)?

%Завдання 11-12 (на звороті): написати функцію для розв'язання задачі. Оцінюється економність обчислень та якість коду.

%Тематика задач: використання циклів, програмування з використанням рекурентних співвідношень,
%обробка списків та рядків.

У завданнях 1-10 вказати, що буде виведено за виконання.

Повідомлення інтерпретатора про помилку позначати ERROR

Завдання 11 та 12 на звороті.

\begin{minipage}[t]{7.5cm}

\begin{center}
\textbf{Завдання 1}
\end{center}

\begin{verbatim}
a, b, c = 7, 8, 9
def f(a, b, c=6):
    print(a, b, c, end=" ")

f(0, 5, a=3)
print(a, b, c)
\end{verbatim}



\begin{center}
\textbf{Завдання 2}
\end{center}

\begin{verbatim}
a,b,c=1,4,6
def h(b):
    a=3
    b-=3
    c=1
    return a+b+c

a,b,c=1,9,7
print(h(b),a,b,c)
\end{verbatim}



\begin{center}
\textbf{Завдання 3}
\end{center}

\begin{verbatim}
def f():
    try:
        res = int("a5")
        return 41
    except Exception: 
        return 1
    except ValueError: 
        return 7
    finally: return 20
    return res

print(f())
\end{verbatim}



\begin{center}
\textbf{Завдання 4}
\end{center}

\begin{verbatim}
n=2
while n<=10:
    n+=2
print(n)
\end{verbatim}

\end{minipage}
\vline \hspace{6mm}
\begin{minipage}[t]{10.5cm}

\begin{center}
\textbf{Завдання 5}
\end{center}

\begin{verbatim}
for b in range(6,6+3,3):
    if b>6:
        continue
    print(b, end=' ')
    b=-3
    if b>=6:
        break
else:
    print(b, end=' ')
print(b, end=' ')
\end{verbatim}



\begin{center}
\textbf{Завдання 6}
\end{center}

\begin{verbatim}
def f(s):
    s1=s
    for i in range(len(s)):
        s1[i]+=1
    return s1

s=[1,2,3]
f(s)
print(s)
\end{verbatim}



\begin{center}
\textbf{Завдання 7}
\end{center}

\begin{verbatim}
h = ('a','b','c','d','e',0,1,2,3)
print(h[-5])
\end{verbatim}



\begin{center}
\textbf{Завдання 8}
\end{center}

\begin{verbatim}
g = ('0','1','2','3','4','5','6','7','8','9')
print(g[-3:10:-2])
\end{verbatim}



\begin{center}
\textbf{Завдання 9}
\end{center}

\begin{verbatim}
h=[10,11,12,13]
h.insert(-3, '10')
print(h)
\end{verbatim}



\begin{center}
\textbf{Завдання 10}
\end{center}
Що буде надруковано за виконання?\par
\begin{verbatim}
h=['0','1','2','3']
del h[:-9]
print(h)
\end{verbatim}

\end{minipage}

%\end {large}
}
\newpage

{{ 27.11.2025 Контрольна робота \No 2, варіант  \No \textbf { 40 }\par}}
{
%\begin {large}

{%Завдання 1-10: Що буде виведено за виконання (повідомлення інтерпретатора про помилку %позначати error)?

%Завдання 11-12 (на звороті): написати функцію для розв'язання задачі. Оцінюється економність обчислень та якість коду.

%Тематика задач: використання циклів, програмування з використанням рекурентних співвідношень,
%обробка списків та рядків.

У завданнях 1-10 вказати, що буде виведено за виконання.

Повідомлення інтерпретатора про помилку позначати ERROR

Завдання 11 та 12 на звороті.

\begin{minipage}[t]{7.5cm}

\begin{center}
\textbf{Завдання 1}
\end{center}

\begin{verbatim}
a, b, c = 9, 8, 7
def g(a, b=6, c):
    print(a, b, c, end=" ")

g(a=2, 3, b=0)
print(a, b, c)
\end{verbatim}



\begin{center}
\textbf{Завдання 2}
\end{center}

\begin{verbatim}
a,b,c=3,4,5
def f(a):
    global c
    a=2
    b+=4
    c=5
    return a+b+c

a,b,c=8,0,2
print(f(a),a,b,c)
\end{verbatim}



\begin{center}
\textbf{Завдання 3}
\end{center}

\begin{verbatim}
def f():
    try:
        res = int(3//0.0)
    except Exception: 
        return 2
    except ValueError: 
        return 6
    else: return 34
    return res

print(f())
\end{verbatim}



\begin{center}
\textbf{Завдання 4}
\end{center}

\begin{verbatim}
m=5
while m<12:
    m+=1
print(m)
\end{verbatim}

\end{minipage}
\vline \hspace{6mm}
\begin{minipage}[t]{10.5cm}

\begin{center}
\textbf{Завдання 5}
\end{center}

\begin{verbatim}
for c in range(-9,-9+5,-2):
    if c>=3:
        continue
    print(c, end=' ')
    c=-6
else:
    print(c, end=' ')
print(c, end=' ')
\end{verbatim}



\begin{center}
\textbf{Завдання 6}
\end{center}

\begin{verbatim}
def f(n):
    if n > 9:
        f(n // 100)
    else:
        print(n % 10, end=' ')

f(987654)
\end{verbatim}



\begin{center}
\textbf{Завдання 7}
\end{center}

\begin{verbatim}
e = ('a','b','c','d','e',0,1,2)
print(e[8])
\end{verbatim}



\begin{center}
\textbf{Завдання 8}
\end{center}

\begin{verbatim}
a = '0123456789'
print(a[13::3])
\end{verbatim}



\begin{center}
\textbf{Завдання 9}
\end{center}

\begin{verbatim}
e=['0','1','2','3']
e.append('11')
print(e)
\end{verbatim}



\begin{center}
\textbf{Завдання 10}
\end{center}
Що буде надруковано за виконання?\par
\begin{verbatim}
g=[10,11,12,13]
g[7:-3]='11'
print(g)
\end{verbatim}

\end{minipage}

%\end {large}
}
\newpage

%end
\end{document}

